% Logic_Functional_Equation.tex
%
% "The Law of Logic as Functional Equation:
%  A Logical Formalization of the Recognition Composition Law."
%
% Draft, April 2026.
%
% Style: MDPI-friendly single-column.

\documentclass[11pt,a4paper]{article}

\usepackage[T1]{fontenc}
\usepackage{lmodern}
\usepackage{amsmath, amssymb, amsthm}
\usepackage[margin=1in]{geometry}
\usepackage{hyperref}
\usepackage{microtype}
\usepackage{enumitem}

\hypersetup{colorlinks=true, linkcolor=black, citecolor=blue, urlcolor=blue}

\newtheorem{theorem}{Theorem}
\newtheorem{lemma}[theorem]{Lemma}
\newtheorem{proposition}[theorem]{Proposition}
\newtheorem{corollary}[theorem]{Corollary}
\newtheorem{conjecture}[theorem]{Conjecture}

\theoremstyle{definition}
\newtheorem{definition}[theorem]{Definition}

\theoremstyle{remark}
\newtheorem{remark}[theorem]{Remark}
\newtheorem{example}[theorem]{Example}

\title{A Functional-Equation Encoding of Logical Consistency on Continuous
Positive-Ratio Comparisons:\\
A Conditional Rigidity Theorem}

\author{Jonathan Washburn\thanks{Recognition Physics Institute, Austin, TX 78701, USA.\\
\texttt{jon@recognitionphysics.org}.}}

\date{April 2026}

\begin{document}

\maketitle

\begin{abstract}
On the continuous-ratio comparison domain, the same forced
functional-equation structure answers two questions. The first is a
cost-physics rigidity question: what cost function on positive ratios
is uniquely determined by consistent composition? The d'Alembert
Inevitability Theorem~\cite{WashburnZlatanovicAllahyarov2026} and the
Canonical Reciprocal Cost Uniqueness
Theorem~\cite{WashburnZlatanovic2026} answer it. The second is a
logic-of-comparison question: what cost function is uniquely
determined by an operational encoding of the four classical
Aristotelian conditions on a continuous comparison operator?

We give the encoding (identity as zero self-cost, non-contradiction as
symmetric single-valued comparison, totality as definiteness on the
open positive quadrant, composition consistency as route-independence
of the combiner), carry two named regularity hypotheses alongside it
(continuity of the operator and polynomial degree-two closure of the
combiner), add a bridge condition (scale invariance), and prove four
translation lemmas reducing the encoded conditions to the hypotheses
of the d'Alembert Inevitability Theorem. The derived cost function
$F(r) := C(r,1)$ then satisfies the Recognition Composition Law
\[
F(xy) + F\!\left(\tfrac{x}{y}\right) = 2F(x) + 2F(y) + c\,F(x)F(y),
\qquad c \in \mathbb{R}.
\]
Under convexity, non-negativity, and the unit log-curvature calibration
$\kappa(F) := \lim_{t \to 0} 2F(e^t)/t^2 = 1$, the family splits into a
bilinear branch ($c \neq 0$) and an additive branch ($c = 0$). The
bilinear branch admits the $\alpha$-family $F_\alpha(x) = (1/\alpha^2)
(\cosh(\alpha \ln x) - 1)$ with $c = 2\alpha^2$ and $\alpha \geq 1$;
its $\alpha = 1$ representative is the canonical reciprocal cost
$J(x) = \tfrac{1}{2}(x + x^{-1}) - 1$. The additive branch admits the
quadratic-logarithm cost $\tfrac{1}{2}(\ln x)^2$. Both branches give
the same calibrated representatives the cost-physics rigidity theorems
produce.

A canonicality theorem (Theorem~\ref{thm:canonicality}) shows that
under the magnitude-of-mismatch interpretation of the comparison
operator, the encoding (L1)--(L4) is the unique reading of the four
classical Aristotelian conditions that preserves their propositional
content. A direct theorem
(Theorem~\ref{thm:rcl-direct}) gives a short self-contained proof
that an operative positive-ratio comparison (continuous, nontrivial,
with $C(x,x) = 0$, $C(x,y) = C(y,x)$, and scale invariance) admitting
a symmetric polynomial combiner of degree at most two satisfies the
Recognition Composition Law. The polynomial restriction is part of
the theorem's hypothesis, not a regularity carried for technical
convenience: Proposition~\ref{prop:finite-closure-necessary} exhibits
a continuous operative comparison ($C(x,y) = (\ln(x/y))^4$) whose
continuous combining rule $\Phi(a, b) = 2a + 2b + 12\sqrt{ab}$ is
closed under iteration on $[0, \infty)$ but is not polynomial of
degree at most two, and the conclusion of the theorem is false for
this example. The identification reported here is therefore not
``every comparison operator satisfying the four classical conditions
yields the Recognition Composition Law''; it is the narrower and
exact identification \emph{finite pairwise polynomial comparison
semantics of classical invariance imply the Recognition Composition
Law}. A short structural theorem
(Theorem~\ref{thm:reality-satisfies-logic}) shows that any
\emph{reality structure} (a triple of states, a real-valued
comparison, and a relation of compositional inference, satisfying
definite states, single-valued comparison, decidability, and
compositional inference) automatically satisfies the four classical
conditions on its comparison operator. Combining the canonicality
theorem, the direct theorem, the reality theorem, and the calibration
machinery, we obtain
(Corollary~\ref{cor:rcl-logic-reality}) the following identification
on the operative domain (continuous positive-ratio reality structures
with finite pairwise polynomial closure): the recognition composition
law, the law of logic, and the structural form of any reality on this
domain are the same mathematical object. The operative domain
contains every pairwise physical measurement, every continuous
likelihood ratio, and every continuous distance or divergence on a
positive-quantity state space. The arithmetic content of the encoding
is recovered in the first companion
paper~\cite{Washburn2026Arithmetic}; the realization-invariance of
that arithmetic across the discrete propositional, modular, ordered,
and categorical realizations is the meta-theorem of the second
companion paper~\cite{Washburn2026UniversalForcing}. The present
paper does not extend to modal or higher-order logics. The discussion
section places the result beside the historical formalisations
(Aristotle to univalent foundations), Wigner's question, the Feynman
name-versus-thing distinction, and the open questions in the
philosophy of mathematics the result bears on.
\end{abstract}

\noindent\textbf{Keywords:} laws of logic; functional equation;
d'Alembert equation; Recognition Composition Law; canonical reciprocal cost;
Wigner's puzzle; foundations of logic.

\medskip

\noindent\textbf{MSC 2020:} 03A05, 03B22, 39B22, 39B52, 39B82.

\section{Introduction}

For nearly two and a half thousand years, the laws of logic have been
treated as a separate object of study from the laws governing the physical
world. Aristotle's three laws of thought, namely the law of identity
($A = A$), the law of non-contradiction ($\neg(A \wedge \neg A)$), and
the law of excluded middle ($A \vee \neg A$), were stated in the
\emph{Prior Analytics} c.~350 BCE \cite{Aristotle}, and have been considered, together
with composition consistency in various forms, the foundational framework
within which valid reasoning occurs.

Successive generations have given these laws increasingly precise formal
articulations. The Stoics, beginning with Chrysippus c.~280 BCE, developed a
propositional calculus around the conditional. Boole's \emph{Laws of
Thought} \cite{Boole1854} recast logic as an algebra. Frege's
\emph{Begriffsschrift} \cite{Frege1879} introduced the predicate calculus
and the quantifiers, providing the notation in which the modern
axiomatization of mathematics is written. Russell and Whitehead's
\emph{Principia Mathematica} \cite{RussellWhitehead1913} established the
logicist program and produced a derivation of arithmetic from the formal
calculus. Hilbert and his school in the 1900s and 1920s axiomatized the
metalanguage of formal proof~\cite{Hilbert1925}. G\"odel's incompleteness
theorems \cite{Goedel1931} established the limits any such axiomatization
can attain. Tarski's semantic conception of truth \cite{Tarski1933} gave the
relationship between formal language and model a precise account. Gentzen
introduced natural deduction and sequent calculus \cite{Gentzen1935},
making the structural content of inference explicit. The Curry--Howard
correspondence \cite{Curry1934, Howard1969} identified proofs with programs.
Lawvere's categorical logic \cite{Lawvere1964} presented logic as a
structure on a category. Girard's linear logic \cite{Girard1987} introduced
resource-aware deduction. The univalent foundations program
\cite{HoTT2013} has explored type-theoretic logic with homotopical content.

Each of these is a substantial formal articulation of one or more aspects of
the laws of logic. Each presents the laws as governing rules within a
designated formal system: a syntax, an algebra, a calculus, a category, a
type theory. None presents the laws of logic as a single unique mathematical
object whose form is forced rather than chosen. The rules describe how
logic operates inside a chosen framework. They do not specify, as a unique
mathematical object, what the laws themselves are.

Feynman, in his Cornell Messenger Lectures of 1964, described the
distinction his father had pressed on him during walks in the woods:
between knowing the name of a bird in five languages and knowing what
the bird is~\cite{Feynman1964}. The prior formalizations of logic give
the laws increasingly precise and increasingly expressive names. Each
remains correct as far as it reaches; the names are mutually consistent
on overlapping domains; the body of knowledge built on each is real and
in some cases indispensable to twentieth-century mathematics. What none
of the formalizations does is exhibit, as a single mathematical
object, what the laws of logic structurally are. The present paper, on
the restricted domain of continuous comparisons of positive ratios,
identifies such an object: the unique continuous cost function
compatible with the encoded Aristotelian conditions plus the named
regularity hypotheses. The identification is conditional on the
encoding; it is mathematically rigorous only on the restricted domain;
it does not generalize automatically to discrete propositional logic or
to other formal systems where the encoding choices may be different.
Within these limits, however, the identification has the form a Feynman
``what the bird is'' answer would have to take: a unique mathematical
object, characterized by structural conditions, with no parameter free
to set after the conditions are fixed.

In this paper we present a formal articulation of a different kind. We
propose a specific operational encoding of four Aristotelian-style
conditions on a continuous comparison operator on positive ratios, carry
through two named regularity hypotheses (continuity of the operator and
polynomial closure of the route-independence combiner), and prove that
under the encoding the conditions force the Recognition Composition Law
$F(xy) + F(x/y) = 2F(x) + 2F(y) + c\,F(x)F(y)$ on the derived cost
function. After branch selection (bilinear $c \neq 0$) and a
multiplicative-coordinate normalization, the unique solution is the
canonical reciprocal cost $J(x) = \tfrac{1}{2}(x + x^{-1}) - 1$. The
parallel additive branch ($c = 0$) yields the quadratic-logarithm cost
$\tfrac{1}{2}(\ln x)^2$ as its unique calibrated solution.

The contribution is a structural \emph{translation theorem}, not a
philosophical reduction. The Aristotelian labels we attach to the
operator-level conditions (Identity, Non-contradiction, Totality,
Composition consistency) are intended as operational analogues on the
positive-ratio comparison setting; we defend each encoding in
Section~\ref{sec:conditions} and acknowledge in
Section~\ref{subsec:not-claimed} what we are \emph{not} claiming. We do
not claim a derivation of classical logic from cost-physics, nor a
derivation of cost-physics from classical logic. The claim is narrower
and structural: under a specific encoding, the same continuous functional
equation is what the conditions admit.

\subsection{Two derivations, one forced combiner family}

Two independently posed rigidity questions, on the continuous-ratio
comparison domain, characterize the same forced combiner family with
the same calibrated branch representatives.

The first is a cost-physics question. Under what minimal hypotheses is
a positive ratio-based cost function uniquely determined by its
composition behaviour? The functional form $F(xy) + F(x/y) = \Phi(F(x),
F(y))$ did not enter the literature as a formalization of logic. It
arose in the analysis of consistent combining rules for ratio-based
costs, in the framework now called Recognition
Geometry~\cite{PardoGuerraThapaWashburnAllahyarov2026,
WashburnZlatanovicAllahyarov2026, WashburnZlatanovic2026}. Two recent
rigidity theorems answer it. The d'Alembert Inevitability
Theorem~\cite{WashburnZlatanovicAllahyarov2026} shows that for
continuous nonconstant $F : \mathbb{R}_{>0} \to \mathbb{R}$ with
$F(1) = 0$ satisfying $F(xy) + F(x/y) = P(F(x), F(y))$ with $P$ a
symmetric polynomial of total degree at most two, the combiner is
forced into $P(u, v) = 2u + 2v + cuv$ for some $c \in \mathbb{R}$.
Higher-degree polynomial combiners are excluded by an explicit
degree-mismatch argument. The Canonical Reciprocal Cost Uniqueness
Theorem~\cite{WashburnZlatanovic2026} shows that under the canonical
$c = 2$ choice with continuity, reciprocal symmetry $F(x) = F(x^{-1})$,
normalization $F(1) = 0$, and the unit log-curvature calibration
$\kappa(F) := \lim_{t \to 0} 2F(e^t)/t^2 = 1$, the unique solution is
$J(x) = \tfrac{1}{2}(x + x^{-1}) - 1$. A companion
result~\cite{WashburnRahnamai2026} reaches the same canonical form in
the ratio-matching optimization setting.

The second question is the one this paper poses. Under what minimal
hypotheses is a positive ratio-based cost function uniquely determined
by an operational encoding of the four classical Aristotelian
conditions on a continuous comparison operator? The encoding is
$C(x, x) = 0$ for identity, $C(x, y) = C(y, x)$ for non-contradiction
read as single-valuedness under reordering, totality of $C$ on the
open positive quadrant for excluded middle in its continuous-domain
form, and the existence of a combining rule $\Phi$ with
$F(xy) + F(x/y) = \Phi(F(x), F(y))$ for composition consistency. These
four conditions, with two named regularity hypotheses (continuity of
the operator and polynomial degree-two closure of $\Phi$) and a bridge
condition (scale invariance), reduce to the hypotheses of the
d'Alembert Inevitability Theorem on the derived cost. The reduction is
the \emph{Translation Theorem} (Theorem~\ref{thm:translation}). The
verification is short and is the technical core of the paper.

The d'Alembert Inevitability Theorem then forces the Recognition
Composition Law family. The Canonical Reciprocal Cost Uniqueness
Theorem then forces $J$ on the bilinear branch under the calibration
and the multiplicative-coordinate normalization $\alpha = 1$. The two
questions, posed for different reasons, characterize the same
mathematical object.

The convergence is at the level of the underlying hypotheses, not
between two derivations from disjoint first principles. The
Aristotelian-style conditions, encoded as we encode them, generate the
same set of hypotheses the cost-physics rigidity theorems use; the
agreement of conclusions is the agreement of inputs, not of two
independent inputs. The point of the paper is therefore narrower than
``logic and cost-physics agree'': it is that one natural reading of
the cost-physics hypotheses on a continuous comparison operator is in
classical logical vocabulary. We discuss this in
Section~\ref{sec:wigner}.

The translation theorem itself is formalised in the canonical
Lean~4 Recognition Science library: the module
\texttt{RecognitionScience.Foundation.LogicAsFunctionalEquation}
contains the full translation from the encoded conditions to the
hypotheses of the d'Alembert Inevitability Theorem, with no residual
\texttt{sorry}s, building against the existing peer-reviewed
\texttt{Inevitability} and \texttt{Cost.FunctionalEquation} modules.
The continuous-combiner extension is formalised in
\texttt{Foundation.GeneralizedDAlembert}, also with no residual
\texttt{sorry}s. The H-side classical input
(\texttt{aczel\_kannappan\_continuous\_dAlembert}) is now a discharged
theorem rather than an axiom; the combiner-side analytic content is
exposed as the explicit hypothesis package
\texttt{ContinuousCombinerAnalysisInputs}, with the same module
formally recording the quartic-log obstruction that prevents the
package from being derived from continuity alone. The earlier
\texttt{Foundation.PolynomialityFromLogic} module, which reduced
polynomiality to a closure-under-iteration analyticity-propagation
class assumption, is superseded by
Proposition~\ref{prop:finite-closure-necessary}, which refutes the
implication that closure under iteration alone forces analyticity. The
analytic reparameterisation counterexample
(Proposition~\ref{prop:analytic-reparam-counterexample}) further shows
that assuming real-analyticity at the origin directly is still not
enough to force polynomial degree at most two. The Lean formalisation
therefore verifies the structural mathematics of the paper; what it
does not and cannot verify is the faithfulness of the Aristotelian
reading of the encoded conditions, which is a choice of definition and
not a theorem about Aristotle. The full module status, including the
H-side discharge and the contents of the hypothesis package, is
reported in Section~\ref{sec:lean-status}.

\subsection{What we are not claiming}\label{subsec:not-claimed}

The result is a translation theorem on a restricted domain, not a
foundational reduction. Four explicit non-claims:

\begin{enumerate}[label=\arabic*., leftmargin=*]
\item Classical logic is not reduced to the cost framework. Valid
reasoning on arbitrary propositions in arbitrary formal systems
remains the province of mathematical logic and is unaffected by the
present result.
\item Finite pairwise polynomial closure of the route-independence
combiner is not asserted as a content of the four classical laws. It
is part of the theorem's hypothesis. Without it, the conclusion of
Theorem~\ref{thm:rcl-direct} is false, witnessed by the operator
$C(x, y) = (\ln(x/y))^4$ with continuous combiner
$\Phi(a, b) = 2a + 2b + 12\sqrt{ab}$
(Proposition~\ref{prop:finite-closure-necessary}). Real-analyticity at
the origin does not repair the gap either
(Proposition~\ref{prop:analytic-reparam-counterexample}). The
identification reported here is therefore not ``every comparison
operator satisfying the four classical conditions yields the
Recognition Composition Law''; it is the narrower and exact
identification ``finite pairwise polynomial comparison semantics of
classical invariance imply the Recognition Composition Law''
(\S\ref{subsec:scope-statement}).
\item The encoding of the Aristotelian conditions is shown to be
canonical under a specific named interpretation of the comparison
operator (Theorem~\ref{thm:canonicality}: under the
magnitude-of-mismatch interpretation, the encoding (L1)--(L4) is the
unique reading that preserves the propositional content of the four
classical laws). The conditionality remaining at this point is on the
choice of interpretation: a non-cost-based reading of the comparison
operator (a directed revision cost, an asymmetric divergence) would
yield a different encoding. Under the standard cost-based reading,
the encoding is determined.
\item The unit log-curvature calibration $\kappa(F) = 1$, with
convexity and non-negativity, does not by itself select $J$. The
calibration leaves a multiplicative-coordinate scaling free, and the
additive branch survives with its own canonical form
$\tfrac{1}{2}(\ln x)^2$. Singling out $J$ requires both branch
selection ($c \neq 0$) and a coordinate normalization fixing the
exponent scale.
\end{enumerate}

What the result does claim is short. Under the encoding plus the
named regularity hypotheses, the Recognition Composition Law is
forced as the unique route-independence combiner family. Its
calibrated solutions on each branch are the canonical forms: $J$ on
the bilinear branch after the $\alpha = 1$ normalization, and
$\tfrac{1}{2}(\ln x)^2$ on the additive branch. The extension to
other carriers (discrete propositional, modular, ordered, categorical)
is the subject of the second companion
paper~\cite{Washburn2026UniversalForcing}, where the
realization-invariance question is given a precise mathematical form.

\subsection{The exact scope of the theorem}\label{subsec:scope-statement}

A point of method that should be stated once, sharply, before any
discussion of implications.

\paragraph{Without finite pairwise polynomial closure, the theorem is
false.} The quartic-log counterexample of
Proposition~\ref{prop:finite-closure-necessary} is an operative
positive-ratio comparison satisfying (L1)--(L4), continuity (R1),
scale invariance (B1), and non-triviality (B2), with a continuous
symmetric combining rule $\Phi(a, b) = 2a + 2b + 12\sqrt{ab}$ that is
closed under iteration on $[0, \infty)$ and is \emph{not} a member of
the Recognition Composition Law family. The analytic-reparameterisation
counterexample of Proposition~\ref{prop:analytic-reparam-counterexample}
shows that even strengthening continuity to real-analyticity at the
origin does not close the gap. Polynomial closure is doing real work,
and we say so explicitly: dropping (R2) makes the conclusion of
Theorem~\ref{thm:translation} fail, by witness, in the same paper.

\paragraph{Why this is not a weakness of the theorem.} The theorem is
not the universal claim ``every comparison operator satisfying the
four classical conditions yields the Recognition Composition Law.''
It is the structurally narrower claim
\begin{quote}
\itshape
finite pairwise polynomial comparison semantics of classical
invariance imply the Recognition Composition Law.
\end{quote}
Finite pairwise polynomial closure is part of the theorem's hypothesis,
not a regularity condition carried for technical convenience. It is
the structural content that distinguishes a finite-algebra comparison
semantics from broader continuous compositional algebras of
positive-ratio comparison. Without it, the operator can still satisfy
the four Aristotelian conditions and still admit a route-independent
combiner; what changes is which combiner family the comparison lives
in, and a non-RCL family is, on the quartic-log example, a perfectly
admissible alternative.

\paragraph{What this clarifies.} The theorem is sharp on its actual
domain: \emph{within the class of comparison operators whose
combiner is a polynomial of total degree at most two}, the
Aristotelian-style encoding forces the Recognition Composition Law
family. Outside that class the conclusion does not hold, and this
paper exhibits the witnesses. The right reading is therefore not
that polynomial closure is an embarrassment to be apologised for, but
that it names the precise compositional class for which the
identification reported here holds. Whether other compositional
classes admit their own rigidity theorems is a separate open question;
this paper does not address it.

\subsection{Companion installments}\label{subsec:companions}

Two companion papers extend the present rigidity theorem in directions
flagged here as out of scope.

The first companion paper~\cite{Washburn2026Arithmetic} takes the
present continuous positive-ratio realization as foundation and
recovers the natural-number tower (Peano arithmetic with order, then
the integers, rationals, reals, and complex numbers) as an internal
construction inside the operative domain, rather than carried in by an
external set-theoretic commitment. The non-claim above that the present
paper does not derive arithmetic is therefore narrower than it reads: it
does not derive arithmetic \emph{within this paper}; arithmetic is the
subject of the next installment, where it is recovered from the same
encoding extended by the orbit construction on the canonical generator.

The second companion paper~\cite{Washburn2026UniversalForcing} states
and proves the Universal Forcing Meta-Theorem: in every admissible
realization of the Law of Logic, the same arithmetic structure is
canonically forced, unique up to canonical isomorphism. The present
paper's continuous positive-ratio realization is the first
mathematically dense instance of that meta-theorem; the discrete
propositional case flagged above as remaining open is treated there as
one further realization of the same abstract framework, alongside
modular, ordered, categorical, and motivating extra-mathematical
realizations. The meta-theorem turns the question of whether the
identification reported here generalises beyond continuous positive
ratios into a precise mathematical statement that can be proved or
refuted realization by realization.

The three papers therefore form a triple. The present paper supplies
the rigidity theorem on continuous positive ratios. The first
companion paper carries that rigidity into the recovery of the number
tower. The second companion paper abstracts the realization framework
and proves invariance of the recovered arithmetic across all
admissible realizations of the Law of Logic.

\subsection{Wigner's question, framed honestly}

In 1960, Eugene Wigner asked why mathematics is ``unreasonably
effective'' in the natural sciences~\cite{Wigner1960}. The puzzle is
that abstract mathematical activity, conducted in pursuit of internal
coherence, gives quantitative descriptions of physical phenomena it
was not designed to describe.

The present result bears on Wigner's question in a narrow way. Under
the encoding adopted here, the Aristotelian-labelled hypotheses on a
comparison operator are the same hypotheses the cost-uniqueness
derivations of \cite{WashburnZlatanovicAllahyarov2026,
WashburnZlatanovic2026} use. The agreement between the two readings of
$J$, one in physics-and-information-geometry vocabulary and one in
Aristotelian-style vocabulary, is not a coincidence between
independent derivations. It is the same hypothesis class characterized
once, with two vocabularies attached.

The paper is therefore a rigidity theorem with an interpretive layer.
The rigidity theorem is the technical core. The interpretive layer
makes the encoding question substantive. Whether the convergence
constitutes a partial dissolution of Wigner's puzzle or only an
observation that the same hypotheses admit two readings on this
domain is a question we take up in Section~\ref{sec:wigner} without
overclaiming either way.

\subsection{Organization}

Section~\ref{sec:operator} introduces the comparison operator and the
derived cost function. Section~\ref{sec:conditions} presents the encoding
of the four Aristotelian-style conditions on the operator, separates the
logical-content conditions (L1)--(L4) from the named regularity
hypotheses (R1) and (R2), and defends each encoding.
Section~\ref{sec:translation} proves the four translation lemmas
connecting the encoded conditions to the hypotheses of the d'Alembert
Inevitability Theorem. Section~\ref{sec:main} states and proves the main
theorem: the encoded conditions force the Recognition Composition Law
family on the derived cost function. Section~\ref{sec:cost} states the
cost corollary, identifying $J$ as the unique calibrated solution on the
bilinear branch after a multiplicative-coordinate normalization, with
the additive-branch counterpart $\tfrac{1}{2}(\ln x)^2$ also displayed.
Section~\ref{sec:finite-pairwise} gives a direct self-contained proof
of the Recognition Composition Law from the operative comparison
conditions plus finite pairwise polynomial closure
(Theorem~\ref{thm:rcl-direct}), and exhibits the quartic-log
counterexample (Proposition~\ref{prop:finite-closure-necessary}) that
shows the polynomial restriction is essential.
Section~\ref{sec:operative} states the operative-domain identification
(Corollary~\ref{cor:rcl-logic-reality}).
Section~\ref{sec:wigner} discusses the historical separation of logic
from physics, the Feynman ``name versus thing'' distinction (with
attention to its dual edge), and the Wigner question with explicit
acknowledgement of the encoding-circularity concern.
Section~\ref{sec:scope} is an account of scope and open questions,
including the final status of finite pairwise polynomial closure after
the closure-under-iteration and analyticity repair attempts both fail,
and is followed by concluding remarks.

\section{The Comparison Operator and Derived Cost}\label{sec:operator}

\begin{definition}[Comparison operator]
\label{def:comparison-operator}
A \emph{comparison operator} is a function
$
C : \mathbb{R}_{>0} \times \mathbb{R}_{>0} \longrightarrow \mathbb{R}.
$
The value $C(x, y)$ is interpreted as the cost (real-valued, possibly
signed) of comparing the two positive quantities $x$ and $y$.
\end{definition}

The two-argument signature reflects the natural data of a comparison: two
quantities being placed into relation. We will pass to a one-argument cost
on ratios via the bridge condition introduced below.

\begin{definition}[Derived cost]
\label{def:derived-cost}
The \emph{derived cost} associated to a comparison operator $C$ is the
function
\[
F : \mathbb{R}_{>0} \longrightarrow \mathbb{R}, \qquad
F(r) := C(r, 1).
\]
\end{definition}

The derived cost fixes the second argument at the multiplicative identity.
Under the bridge condition (scale invariance) introduced in
Section~\ref{sec:conditions}, the derived cost depends only on the ratio of
the two arguments, and the two-argument operator can be reconstructed
from $F$ via $C(x, y) = F(x/y)$.

\section{The Aristotelian Conditions, Encoded}\label{sec:conditions}

The encoding of the four Aristotelian-style conditions on a continuous
comparison operator on positive ratios has two halves. The
\emph{logical content} is four conditions that carry the structural
import of the classical laws on this domain. The \emph{regularity
hypotheses} are two analytical conditions (continuity and polynomial
degree-two closure of the route-independence combiner) carried
alongside the logical content. The regularity hypotheses are
\emph{not} claimed to be contents of the classical laws.

The split tells the reader which parts of the rigidity theorem do the
structural work attributed to the classical laws and which parts are
regularity demands carried by hand.

\subsection{Logical content versus regularity hypotheses}\label{subsec:logical-vs-regularity}

\paragraph{Logical content.} On a continuous comparison operator on
positive ratios, the four classical Aristotelian conditions translate
as follows. Each translation is operational, not derivational.
Section~\ref{subsec:encoding-defenses} defends each one.

\begin{enumerate}[label=(L\arabic*), leftmargin=*]
\item \textbf{Identity, encoded as zero self-cost}: $C(x, x) = 0$.
The cost of comparing a quantity to itself is zero.
\item \textbf{Non-contradiction, encoded as symmetric single-valued
comparison}: $C(x, y) = C(y, x)$. The cost of a comparison does not
depend on which of the two quantities is named first; when read as a
magnitude of mismatch (rather than as a directed revision cost), the
comparison is single-valued under reordering of arguments. We use this
longer label throughout the paper to flag that the identification with
non-contradiction is an operational encoding, not a claim of literal
equivalence to the propositional law $\neg(A \wedge \neg A)$.
\item \textbf{Totality, encoded as definiteness on the open
quadrant}: $C$ is defined on the entire open positive quadrant
$\mathbb{R}_{>0} \times \mathbb{R}_{>0}$ and takes real values. Every
well-formed comparison returns a definite real value; there are no
input pairs on which the comparison fails to be defined or returns an
``undefined'' marker. This is the continuous-domain encoding of
excluded middle: on every input, the comparison takes a value (the
analogue of ``$A$ holds or $\neg A$ holds''), with no third
``undefined'' possibility. We note that, under the type signature of
Definition~\ref{def:comparison-operator} (a function on the open
positive quadrant), this condition is \emph{automatic}: the content
the Aristotelian labelling is attached to is the choice of the
definition itself, rather than a separate axiom. The encoding of
excluded middle is therefore the structure of the operator's signature
and codomain, not a theorem about the operator.
\item \textbf{Composition consistency (route-independence)}: there is a
combining rule $\Phi$ such that the cost of composite comparisons in the
forward-and-backward symmetric form satisfies
$F(xy) + F(x/y) = \Phi(F(x), F(y))$ on the derived cost. The cost of a
composite comparison is a function of its constituent costs alone; the
route by which the composite is assembled does not contribute additional
cost.
\end{enumerate}

\paragraph{Regularity hypotheses.} The following two hypotheses are
analytical conditions, named explicitly, that we carry alongside the
logical content. We do not claim either is forced by the four classical
laws.

\begin{enumerate}[label=(R\arabic*), leftmargin=*]
\item \textbf{Continuity}: the comparison operator $C$ is jointly
continuous on $\mathbb{R}_{>0} \times \mathbb{R}_{>0}$. This is a
regularity demand on the operator, not a logical content. It supplies
the analytic infrastructure on which the d'Alembert-type rigidity machinery
runs.
\item \textbf{Polynomial closure of $\Phi$ of degree at most two}: the
combining rule $\Phi$ in (L4) is a symmetric polynomial of total degree
at most two on $\mathbb{R}^2$. This is a regularity restriction inherited
from the d'Alembert Inevitability Theorem
\cite{WashburnZlatanovicAllahyarov2026}. We carry it explicitly. The
quartic-log counterexample of
Proposition~\ref{prop:finite-closure-necessary} shows that the
restriction is essential: without it, an operative positive-ratio
comparison can have a non-polynomial continuous combiner that escapes
the Recognition Composition Law family. The analytic reparameterisation
counterexample of Proposition~\ref{prop:analytic-reparam-counterexample}
shows that even real-analyticity of the combiner at the origin does not
force polynomial degree at most two. Thus polynomial closure remains an
essential finite-algebra hypothesis, not a removable regularity condition.
\end{enumerate}

\paragraph{Bridge and non-triviality.} Two further conditions, of a
different character from either logical content or regularity, complete
the setup.

\begin{enumerate}[label=(B\arabic*), leftmargin=*]
\item \textbf{Scale invariance}: $C(\lambda x, \lambda y) = C(x, y)$ for
all $\lambda > 0$. This is a structural bridge that lets the
two-argument operator descend consistently to a one-argument cost on the
multiplicative group of positive ratios. It is motivated by the
scale-free character of the classical laws (the laws make no reference
to absolute magnitude) but is not itself one of the classical laws.
\item \textbf{Non-triviality}: there exists $x_0 > 0$ with
$F(x_0) \neq 0$. This excludes the constant-zero operator, which
vacuously satisfies all the other conditions.
\end{enumerate}

The split between (L\,$*$), (R\,$*$), (B\,$*$) is preserved through
the rest of the paper. The main rigidity theorem is conditional on all
seven. Each lemma below names which of the seven it uses, and the
regularity hypotheses are flagged each time they enter.

\subsection{Defenses of the encodings}\label{subsec:encoding-defenses}

The four encodings of the classical laws are open to alternatives. We
defend each below.

\paragraph{(L1) Identity.} The identification of $C(x, x) = 0$ with the
law of identity $A = A$ is the cleanest of the four. The cost of
comparing a quantity to itself is zero. The identification fixes the
normalization of the cost scale; nothing further is at stake.

\paragraph{(L2) Non-contradiction.} The identification of $C(x, y) =
C(y, x)$ with non-contradiction is contestable. Two readings of
``comparison'' need to be distinguished. First, comparison as a
magnitude of mismatch between two configurations: how different are
$x$ and $y$? Under this reading, the cost is symmetric in its arguments;
the question of which configuration is named first is immaterial because
the magnitude of mismatch is a property of the pair, not of the order
of presentation. Second, comparison as a directed revision: what does it
cost to update $x$ to match $y$, or $y$ to match $x$? Under this reading,
the cost is plausibly asymmetric and (L2) does not hold.

We adopt the first reading. This is a choice. The choice aligns with the
Aristotelian formulation in the following sense: if a single comparison
is held to have two distinct values $C(x, y) \neq C(y, x)$, the same
comparison would have two simultaneous valuations, in violation of the
single-valuedness that any well-formed predication requires. The
asymmetric reading does not violate this; it simply says that ``compare
$x$ to $y$'' and ``compare $y$ to $x$'' are different comparisons with
distinct asymmetric values. We do not claim the symmetric reading is the
only natural one. We claim it is the natural one for cost as a magnitude
of mismatch, and we adopt it for that purpose.

\paragraph{(L3) Totality.} The identification of (L3) with excluded
middle requires explicit care. Classical excluded middle says that for
every well-formed proposition $A$, either $A$ or $\neg A$ holds. On a
continuous comparison operator, the natural analogue is that for every
input pair $(x, y)$, the comparison returns a value: there are no input
configurations on which the comparison is undefined.

We \emph{separate this from continuity}. A function can be total,
i.e.\ defined everywhere on its domain, without being continuous. Step functions,
the Dirichlet function, and many other discontinuous functions are total
on their domains. Continuity is a regularity property; totality is the
content we identify with excluded middle.

The formal definitions below therefore separate (L3)
\emph{Totality} (the operator is defined on all of $\mathbb{R}_{>0}^2$,
automatic from Definition~\ref{def:comparison-operator}) from (R1)
\emph{Continuity} (joint continuity on the open quadrant).

\paragraph{(L4) Composition consistency.} The classical content of
composition consistency is the transitivity and well-formedness of
inference: if a composite reasoning step is valid, its validity is
determined by the validity of its constituent steps, with no additional
contribution from the order or association of the assembly. On a
continuous comparison operator, the natural analogue is route-independence
of the composite cost: the cost of comparing the assembled compositions
$xy$ and $x/y$ is a function of the costs of the constituent ratios
$F(x)$ and $F(y)$, not of any further data about how the composition
was formed.

This is what (L4) asserts: the combining rule $\Phi$ exists. The
existence of $\Phi$, taken as a symmetric continuous function on
$\operatorname{Range}(F)^2$, is the content we identify with composition
consistency.

The further restriction that $\Phi$ be a polynomial of total degree at
most two is \emph{not} part of (L4). It is the regularity hypothesis (R2),
named separately. (R2) is what permits the
d'Alembert Inevitability Theorem to apply directly.
Proposition~\ref{prop:finite-closure-necessary} shows that the
polynomial restriction is essential, not redundant. Proposition~\ref{prop:analytic-reparam-counterexample}
shows that real-analyticity of the combiner at the origin also fails to
force polynomial degree at most two. Thus (R2) is the exact finite
pairwise algebra condition used by the theorem.

\paragraph{Summary of the split.} The four classical Aristotelian-style
conditions are (L1)--(L4). The two analytic hypotheses (R1) and (R2)
are named regularity demands. (B1) and (B2) are a structural bridge
and a degeneracy exclusion. The main rigidity theorem
(Section~\ref{sec:main}) uses all seven. The claim is that (L1)--(L4)
augmented by (R1), (R2), (B1), (B2) force the Recognition Composition
Law; (L1)--(L4) alone do not. The encoding (L1)--(L4) is the new
content of the paper; the rigidity then follows from the published
d'Alembert Inevitability Theorem.

\subsection{The conditions, formally}

\begin{definition}[Identity]
\label{def:identity}
A comparison operator $C$ \emph{satisfies Identity} if
\[
C(x, x) = 0 \quad \text{for all } x > 0.
\]
\end{definition}

\begin{definition}[Non-contradiction]
\label{def:noncontradiction}
A comparison operator $C$ \emph{satisfies Non-contradiction} if
\[
C(x, y) = C(y, x) \quad \text{for all } x, y > 0.
\]
\end{definition}

\begin{definition}[Totality (logical content)]
\label{def:totality}
A comparison operator $C$ \emph{satisfies Totality} if it is defined on
all of $\mathbb{R}_{>0} \times \mathbb{R}_{>0}$. (This is automatic from
Definition~\ref{def:comparison-operator}, which defines $C$ as a
function on the open positive quadrant. We name the condition explicitly
because it is the operational content we identify with excluded middle:
every input pair returns a value, with no third ``undefined'' alternative.)
\end{definition}

\begin{definition}[Continuity (regularity hypothesis (R1))]
\label{def:continuity}
A comparison operator $C$ \emph{satisfies Continuity} if the two-variable
map $(x, y) \longmapsto C(x, y)$ is jointly continuous on
$\mathbb{R}_{>0} \times \mathbb{R}_{>0}$. This is a regularity hypothesis,
not a logical content; it is named separately from Totality
(Definition~\ref{def:totality}).
\end{definition}

\begin{definition}[Composition consistency (logical content)]
\label{def:composition-consistency}
A comparison operator $C$ \emph{satisfies Composition consistency} if
there exists a function $\Phi : \mathbb{R}^2 \to \mathbb{R}$ such that
the derived cost $F(r) := C(r, 1)$ satisfies
\[
F(xy) + F\!\left(\tfrac{x}{y}\right) = \Phi(F(x), F(y))
\quad \text{for all } x, y > 0.
\]
The function $\Phi$ is the \emph{combining rule}.
\end{definition}

\begin{definition}[Polynomial closure of $\Phi$ (regularity hypothesis (R2))]
\label{def:route-independence}
A comparison operator $C$ that satisfies Composition consistency
\emph{satisfies Polynomial closure of degree at most two} if the
combining rule $\Phi$ extends to a symmetric polynomial $P \in
\mathbb{R}[u, v]$ of total degree at most two on $\mathbb{R}^2$ for which
\[
F(xy) + F\!\left(\tfrac{x}{y}\right) = P(F(x), F(y))
\quad \text{for all } x, y > 0.
\]
This is a regularity hypothesis on the combining rule, not a content of
Composition consistency. We adopt it in the main theorem because the
d'Alembert Inevitability
Theorem~\cite{WashburnZlatanovicAllahyarov2026} requires polynomiality of
degree at most two; we discuss its status as an assumption in
Section~\ref{subsec:conjecture}.
\end{definition}

\begin{remark}
The split between Composition consistency (Definition~\ref{def:composition-consistency},
existence of a combining rule) and Polynomial closure
(Definition~\ref{def:route-independence}, symmetric polynomial of degree
$\leq 2$) is the central conceptual move of the encoding. The first is
the operational content of route-independence on a continuous comparison
operator. The second is the regularity hypothesis (R2) carried alongside.
The symmetry of $\Phi$ on $\operatorname{Range}(F)^2$ is a consequence of
Non-contradiction at the combiner level (see
\cite[Lemma~1]{WashburnZlatanovicAllahyarov2026}); under (R2) it extends
to the full polynomial $P$ on $\mathbb{R}^2$ by the standard argument
that a polynomial vanishing on a product of nondegenerate intervals
vanishes identically.
\end{remark}

\subsection{The bridge condition and non-triviality}

\begin{definition}[Scale invariance]
\label{def:scale-invariance}
A comparison operator $C$ is \emph{scale-invariant} if
\[
C(\lambda x, \lambda y) = C(x, y) \quad \text{for all } x, y, \lambda > 0.
\]
\end{definition}

The four classical laws of logic make no reference to absolute magnitude.
They are statements about the relations among arguments, not about the
arguments' overall scale. To express them as constraints on the
two-argument operator $C$ in a way that descends consistently to a
one-argument cost on positive ratios, we require that $C$ depend only on
the ratio of its arguments. This is the content of scale invariance: it is
the bridge from the operator-level form of the four conditions to their
form on the multiplicative group $\mathbb{R}_{>0}$. Without scale
invariance, comparisons would be sensitive to absolute magnitude, in
violation of the scale-free character of the classical laws.

\begin{definition}[Non-triviality]
\label{def:nontriviality}
A comparison operator $C$ is \emph{non-trivial} if there exists $x_0 > 0$
with $F(x_0) = C(x_0, 1) \neq 0$.
\end{definition}

\begin{remark}
Without non-triviality, the constant-zero operator $C \equiv 0$ vacuously
satisfies all the constraints. Non-triviality excludes this degenerate
case and is the analogue of the requirement (in
\cite{WashburnZlatanovicAllahyarov2026}) that $F$ be nonconstant.
\end{remark}

\begin{remark}[Non-triviality is now derived, not posited]
\label{rem:nontriviality-derived}
The clause (B2) is no longer carried as a free posit at the foundation
layer. The companion absolute-floor closure
\cite{Washburn2026FloorClosure} reduces it to two named preconditions:
the meta-language admits at least two distinct propositions
\cite{Washburn2026SelfBootstrap}, and the universe of discourse
contains more than one element \cite{Washburn2026Specifiability}. On
every inhabited carrier, these jointly supply the operator-level
non-triviality clause through the Lean bridge structure
\texttt{Foundation.NonTrivialityFromDistinguishability.SatisfiesLawsOfLogicAbsoluteFloor},
which combines an absolute-floor witness on the positive-ratio
carrier with a detection condition that the comparison operator
returns a non-zero value on the witnessed pair. The discharge theorem
is \texttt{existing\_of\_absoluteFloor}.
\end{remark}

\begin{definition}[The encoded conditions on a comparison operator]
\label{def:laws-of-logic}
A comparison operator $C: \mathbb{R}_{>0} \times \mathbb{R}_{>0} \to
\mathbb{R}$ \emph{satisfies the encoded conditions for the laws of logic
on continuous positive-ratio comparisons} if it satisfies:
\begin{itemize}[leftmargin=*]
\item the four logical-content conditions: Identity
(Definition~\ref{def:identity}), Non-contradiction
(Definition~\ref{def:noncontradiction}), Totality
(Definition~\ref{def:totality}), and Composition consistency
(Definition~\ref{def:composition-consistency});
\item the two regularity hypotheses: Continuity
(Definition~\ref{def:continuity}, hypothesis (R1)) and Polynomial
closure of the combiner of degree at most two
(Definition~\ref{def:route-independence}, hypothesis (R2));
\item the bridge condition Scale invariance
(Definition~\ref{def:scale-invariance}, (B1));
\item the non-triviality clause
(Definition~\ref{def:nontriviality}, (B2)).
\end{itemize}
The phrasing ``satisfies the encoded conditions'' is intended to flag
that the result is conditional on the encoding: see
Section~\ref{subsec:encoding-defenses} for the defenses,
Section~\ref{sec:canonicality} for the canonicality theorem that
sharpens the conditionality, and Section~\ref{subsec:not-claimed} for
what is not claimed.
\end{definition}

\section{Canonicality of the Encoding}\label{sec:canonicality}

The encoding presented in Section~\ref{sec:conditions} maps the four
classical Aristotelian conditions to four structural conditions on the
comparison operator. We have so far called the encoding
``operational'' rather than ``derivational.'' The purpose of the
present section is to establish the precise sense in which the
encoding is canonical: under a single named interpretation of the
comparison operator, the encoding is the unique reading of the four
classical conditions that preserves their propositional content.

\subsection{The magnitude-of-mismatch interpretation}

\begin{definition}[Magnitude-of-mismatch interpretation]
\label{def:mismatch-interpretation}
A comparison operator
$C : \mathbb{R}_{>0} \times \mathbb{R}_{>0} \to \mathbb{R}$ is read as
a \emph{magnitude-of-mismatch} if its codomain is interpreted as the
structural distance between the two arguments, with the following
three properties:
\begin{enumerate}[label=(M\arabic*),leftmargin=*]
\item \emph{Trivial value at match.} The trivial element of the
codomain, taken to be $0 \in \mathbb{R}$, encodes ``perfect match''
between the two arguments. Non-trivial values encode progressively
greater divergence from match.
\item \emph{Pair-symmetry under natural exchange.} The mismatch is a
property of the unordered pair $\{x, y\}$ rather than of the ordered
pair $(x, y)$, under the natural exchange symmetry the operator
inherits from the symmetry of ``same magnitude'' between two positive
quantities.
\item \emph{Determinability under composition.} The mismatch of a
composite comparison, assembled by an operation that combines two
constituent comparisons, is determined by the mismatches of its
constituents and the assembly operation.
\end{enumerate}
\end{definition}

\begin{remark}
The magnitude-of-mismatch interpretation is the standard reading of
cost functions in information geometry, statistics, and physics. The
Kullback--Leibler divergence, the Bregman divergences, the
squared-error penalty in least squares, the chi-squared statistic,
and the recognition cost $J$ all fit this interpretation. The reading
is therefore not an ad hoc choice; it is the standard reading of a
cost-based comparison.
\end{remark}

\subsection{The canonicality theorem}

\begin{theorem}[Canonicality of the encoding]
\label{thm:canonicality}
Let
$C : \mathbb{R}_{>0} \times \mathbb{R}_{>0} \to \mathbb{R}$ be a
comparison operator interpreted as a magnitude-of-mismatch in the
sense of Definition~\ref{def:mismatch-interpretation}. The unique
reading of the four classical Aristotelian conditions (identity,
non-contradiction, excluded middle, composition consistency) on $C$
that preserves the propositional content of each condition is the
encoding (L1)--(L4) of Section~\ref{sec:conditions}.
\end{theorem}

\begin{proof}
We argue case by case for each condition.

\paragraph{Identity.} The propositional content of identity is ``a
thing is itself.'' Under (M1), ``is itself'' translates to ``the
mismatch with itself is the trivial element of the codomain.'' The
trivial element is $0$. So the identity condition translates to
$C(x, x) = 0$ for every $x > 0$. Any alternative
($C(x, x) = \delta \neq 0$) would assert that a thing has positive
structural distance from itself, in violation of the propositional
content of identity together with (M1). The encoding (L1) is therefore
the unique reading.

\paragraph{Non-contradiction.} The propositional content of
non-contradiction is ``the same instance cannot simultaneously hold
two distinct values.'' On a two-argument comparison operator, the
most direct way for ``the same instance'' to have two distinct values
is for $C(x, y)$ and $C(y, x)$ to differ. Under (M2), the mismatch is
a property of the unordered pair $\{x, y\}$; thus $C(x, y)$ and
$C(y, x)$ must agree, since they are two readings of the same
unordered-pair mismatch. Any alternative ($C(x, y) \neq C(y, x)$)
would assign two distinct values to the same instance, in violation
of the propositional content. The encoding (L2) is therefore the
unique reading.

\paragraph{Excluded middle.} The propositional content of excluded
middle is ``every well-formed instance has a definite outcome; there
is no third option.'' On a comparison operator with codomain
$\mathbb{R}$, ``definite outcome'' means a single real value. ``No
third option'' means there is no input pair in the domain on which
the operator returns ``undefined'' or some third codomain element.
This translates to totality of $C$ on
$\mathbb{R}_{>0} \times \mathbb{R}_{>0}$. On a continuous comparison
operator, the continuous version of ``no third option'' forbids
tearing of the value across infinitesimal input changes: a
discontinuity in $C$ would correspond to two simultaneous limit
values (left and right) at the discontinuity, in violation of the
propositional content. Continuity is therefore the continuous-domain
refinement of ``no third option.'' The encoding (L3) plus the
regularity hypothesis (R1) together capture the continuous reading;
(L3) alone is the totality content.

\paragraph{Composition consistency.} The propositional content of
composition consistency is the well-formedness of inferential
composition: a composite reasoning step is determinable from its
constituents, with no dependence on the order or association of the
assembly. Under (M3), the mismatch of a composite comparison is
determined by the mismatches of its constituents and the assembly
operation. On the multiplicative group $\mathbb{R}_{>0}$, the natural
forward-and-backward symmetric assembly of two comparisons is the
pair $(xy, x/y)$; the costs of these two assemblies are summed to
give the cost of the composite. The d'Alembert form
$F(xy) + F(x/y) = \Phi(F(x), F(y))$ expresses, on the derived cost
$F$, the determinability of the composite cost from the constituent
costs $F(x)$ and $F(y)$ via a combining rule $\Phi$. Any alternative
form (route-dependent, asymmetric in unintended ways, or with
composite cost depending on data not derivable from $F(x), F(y)$)
would violate either (M3) or the propositional content. The encoding
(L4) is therefore the unique reading.
\end{proof}

\begin{corollary}[The encoding is determined under the standard cost reading]
\label{cor:encoding-determined}
Under the magnitude-of-mismatch interpretation, the encoded conditions
(L1)--(L4) are the unique mapping of the four classical Aristotelian
conditions onto a continuous comparison operator on positive ratios.
The encoding is not one choice among many; it is the unique encoding
determined by the propositional content of the classical conditions
together with the structural content of the magnitude-of-mismatch
interpretation.
\end{corollary}

\subsection{Implications for the conditional-rigidity framing}

The encoding hedge of Section~\ref{subsec:not-claimed} (item 3, which
read that the encoding ``is not asserted to be canonical or unique'')
is now refined by Theorem~\ref{thm:canonicality}. The encoding is
canonical and unique \emph{under the magnitude-of-mismatch
interpretation}. The remaining residual is therefore narrower: the
encoding depends on the choice of interpretation of the comparison
operator, and Theorem~\ref{thm:canonicality} fixes it under the
standard cost-based interpretation. A reading of the comparison
operator as something other than a magnitude of mismatch (a directed
revision cost, an asymmetric divergence in some non-standard sense,
etc.) would yield a different encoding and a different rigidity
result. The standard reading produces the standard encoding (L1)--(L4)
and the rigidity theorems of the present paper.

The strengthened framing for the rest of the paper: under the standard
cost-based reading of the comparison operator, the encoding is
determined, and the result of Theorem~\ref{thm:translation} and
Theorem~\ref{thm:main} is unconditional on the choice of encoding
within that reading. The remaining conditionality is on the regularity
hypotheses (R1) and (R2), and on the magnitude-of-mismatch
interpretation itself, named explicitly above.

\section{The Translation Theorem}\label{sec:translation}

We prove four translation lemmas, then combine them.

\begin{lemma}[Identity implies normalization]
\label{lem:identity-normalization}
If $C$ satisfies Identity, then $F(1) = 0$.
\end{lemma}

\begin{proof}
Apply Identity at $x = 1$: $F(1) = C(1, 1) = 0$.
\end{proof}

\begin{lemma}[Non-contradiction with scale invariance implies reciprocity]
\label{lem:reciprocity}
If $C$ satisfies Non-contradiction and Scale invariance, then $F$ is
reciprocal-symmetric:
\[
F(x) = F(x^{-1}) \quad \text{for all } x > 0.
\]
\end{lemma}

\begin{proof}
Fix $x > 0$. Then
\begin{align*}
F(x)
  &= C(x, 1) && \text{(definition of $F$)}\\
  &= C(1, x) && \text{(Non-contradiction)}\\
  &= C(x^{-1} \cdot 1, \; x^{-1} \cdot x)
                          && \text{(Scale invariance with $\lambda = x^{-1}$)}\\
  &= C(x^{-1}, 1) && \text{(simplification)}\\
  &= F(x^{-1}) && \text{(definition of $F$)}.
\end{align*}
The third step uses Scale invariance applied to the pair $(1, x)$ with the
factor $\lambda = x^{-1} > 0$, taking the form $C(1, x) = C(x^{-1} \cdot 1,
x^{-1} \cdot x)$. The fourth step uses $x^{-1} \cdot 1 = x^{-1}$ and
$x^{-1} \cdot x = 1$.
\end{proof}

\begin{lemma}[Continuity of $C$ implies continuity of $F$]
\label{lem:continuity}
If $C$ satisfies Continuity (R1), then $F$ is continuous on
$\mathbb{R}_{>0}$.
\end{lemma}

\begin{proof}
$F$ is the composition $F = C \circ \iota$, where
$\iota : \mathbb{R}_{>0} \to \mathbb{R}_{>0} \times \mathbb{R}_{>0}$,
$\iota(r) = (r, 1)$. The map $\iota$ is continuous on $\mathbb{R}_{>0}$
(both component functions are continuous, the first being the identity
and the second a constant). The map $C$ is jointly continuous on
$\mathbb{R}_{>0} \times \mathbb{R}_{>0}$ by (R1). The composition of
continuous functions is continuous, so $F$ is continuous on
$\mathbb{R}_{>0}$.
\end{proof}

\begin{lemma}[Composition consistency with polynomial closure transfers to $F$]
\label{lem:route-independence-transfer}
If $C$ satisfies Composition consistency
(Definition~\ref{def:composition-consistency}) together with Polynomial
closure of $\Phi$ of degree at most two
(Definition~\ref{def:route-independence}, hypothesis (R2)), then $F$
satisfies the polynomial composition law
\[
F(xy) + F\!\left(\tfrac{x}{y}\right) = P(F(x), F(y))
\]
for all $x, y > 0$, with $P$ symmetric and of total degree at most two.
\end{lemma}

\begin{proof}
Immediate from Definition~\ref{def:route-independence}: the polynomial
composition law on $F$ is the content of (R2) as stated.
\end{proof}

\begin{theorem}[Translation Theorem]
\label{thm:translation}
Let $C$ be a comparison operator satisfying the encoded conditions
(Definition~\ref{def:laws-of-logic}). Then the derived cost function
$F(r) := C(r, 1)$ satisfies all the hypotheses of the d'Alembert
Inevitability
Theorem~\cite[Theorem~3]{WashburnZlatanovicAllahyarov2026}: $F$ is
continuous on $\mathbb{R}_{>0}$, nonconstant, with $F(1) = 0$, and
satisfies a polynomial composition law $F(xy) + F(x/y) = P(F(x), F(y))$
for some symmetric polynomial $P$ of total degree at most two.
\end{theorem}

\begin{proof}
$F(1) = 0$ by Lemma~\ref{lem:identity-normalization} (Identity).
Continuity of $F$ on $\mathbb{R}_{>0}$ is
Lemma~\ref{lem:continuity} ((R1)). $F$ is nonconstant by Non-triviality
(B2) together with $F(1) = 0$: there exists $x_0 > 0$ with
$F(x_0) \neq 0 = F(1)$, so $F$ takes at least two distinct values. The
polynomial composition law on $F$ with symmetric $P$ of total degree at
most two is Lemma~\ref{lem:route-independence-transfer} (Composition
consistency together with (R2)). All hypotheses of
\cite[Theorem~3]{WashburnZlatanovicAllahyarov2026} are met.

(Reciprocity $F(x) = F(x^{-1})$ from
Lemma~\ref{lem:reciprocity} is also a consequence of the encoded
conditions, but is not required as an additional hypothesis: in
\cite[Lemma~1]{WashburnZlatanovicAllahyarov2026} reciprocity follows from
symmetry of $P$, and is therefore part of the conclusion of the
Inevitability Theorem rather than its hypothesis.)
\end{proof}

\section{Main Theorem: The Recognition Composition Law as the Unique Combiner Family}\label{sec:main}

\begin{theorem}[Translation-and-rigidity theorem]
\label{thm:main}
Let $C$ be a comparison operator satisfying the encoded conditions
(Definition~\ref{def:laws-of-logic}). Then there exists $c \in \mathbb{R}$
such that the derived cost function $F$ satisfies the Recognition
Composition Law
\[
F(xy) + F\!\left(\tfrac{x}{y}\right) = 2F(x) + 2F(y) + c\,F(x)F(y)
\quad \text{for all } x, y > 0. \tag{$\mathrm{RCL}$}
\]
The polynomial $P$ underlying (R2) is uniquely determined as
$P(u, v) = 2u + 2v + c\,uv$.
\end{theorem}

\begin{proof}
By Theorem~\ref{thm:translation}, the derived cost $F$ satisfies the
hypotheses of the d'Alembert Inevitability Theorem
\cite[Theorem~3]{WashburnZlatanovicAllahyarov2026}: $F$ is continuous on
$\mathbb{R}_{>0}$, nonconstant, with $F(1) = 0$, and satisfies a polynomial
composition law $F(xy) + F(x/y) = P(F(x), F(y))$ with $P$ symmetric of total
degree at most two. The Inevitability Theorem then forces
\[
P(u, v) = 2u + 2v + c\,uv \quad \text{for some } c \in \mathbb{R}.
\]
Substituting into the polynomial composition law gives $(\mathrm{RCL})$.

The polynomial $P$ is uniquely determined by $F$: by
\cite[Corollary~5]{WashburnZlatanovicAllahyarov2026}, if two polynomials
$P, P'$ satisfy $F(xy) + F(x/y) = P(F(x), F(y)) = P'(F(x), F(y))$ for all
$x, y > 0$ with $F$ continuous nonconstant, then $P \equiv P'$ on
$\mathbb{R}^2$. Hence the value of $c$ is determined by $F$.
\end{proof}

\begin{remark}[The structural content]
\label{rem:structural-content}
Theorem~\ref{thm:main} says the four encoded conditions on a continuous
comparison operator, with scale invariance and the exclusion of the
trivial operator, admit \emph{exactly one} functional form for the
route-independence combiner, parametrised by a single real number
$c$. The combining rule is forced. Higher-degree polynomial combiners
are excluded by the Inevitability Theorem
\cite[Theorem~1]{WashburnZlatanovicAllahyarov2026}; lower-degree
alternatives fall into the $c = 0$ case
\cite[Remark~4]{WashburnZlatanovicAllahyarov2026}. The parameter $c$
is a curvature degree of freedom, fixed in Section~\ref{sec:cost} by a
calibration condition.
\end{remark}

\begin{remark}[Reciprocity and the role of scale invariance]
\label{rem:reciprocity}
By Lemma~\ref{lem:reciprocity}, the derived cost $F$ of any comparison
operator satisfying the encoded conditions is reciprocal-symmetric:
$F(x) = F(x^{-1})$. The reciprocity follows from Non-contradiction
together with Scale invariance. Non-contradiction alone gives
$C(x, y) = C(y, x)$ (a symmetry of the two-argument operator); the
descent to $F(x) = F(x^{-1})$ requires scale invariance as the bridge
condition that lets the classical laws, scale-free at the level of
their statement, be written as constraints on the multiplicative group
$\mathbb{R}_{>0}$.
\end{remark}

\section{RCL as the Finite Pairwise Algebra of Positive-Ratio Comparison}\label{sec:finite-pairwise}

Theorem~\ref{thm:main} establishes the Recognition Composition Law
family by routing through the d'Alembert Inevitability
Theorem~\cite{WashburnZlatanovicAllahyarov2026}. We now give a direct
self-contained proof from the encoded conditions, using only the
operator-level structure plus a single named regularity hypothesis
(finite pairwise polynomial closure of the combiner). The direct proof
is shorter than the d'Alembert route, and its hypothesis is the precise
structural condition that distinguishes the Recognition Composition Law
from broader continuous-combiner alternatives. We also display a
counterexample (Proposition~\ref{prop:finite-closure-necessary}) showing
that finite pairwise polynomial closure is not redundant: dropping it
admits comparison structures with non-polynomial combiners that are not
in the Recognition Composition Law family.

\subsection{Operative positive-ratio comparison}

\begin{definition}[Operative positive-ratio comparison]
\label{def:operative-comparison}
An \emph{operative positive-ratio comparison} is a continuous,
nontrivial map
\[
C : \mathbb{R}_{>0} \times \mathbb{R}_{>0} \to \mathbb{R}
\]
satisfying
\[
C(x, x) = 0, \qquad C(x, y) = C(y, x), \qquad
C(\lambda x, \lambda y) = C(x, y)
\]
for all $x, y, \lambda > 0$. Its derived cost is
\[
F(r) := C(r, 1).
\]
\end{definition}

The three conditions on $C$ in Definition~\ref{def:operative-comparison}
are exactly Identity (L1), Non-contradiction (L2), and Scale invariance
(B1) of Section~\ref{sec:conditions}, with continuity as the
regularity hypothesis (R1) and totality of $C$ as Definition~\ref{def:totality}.
The definition isolates the structural minimum we need for the direct
proof below.

\begin{definition}[Finite pairwise polynomial closure]
\label{def:finite-pairwise-closure}
An operative positive-ratio comparison $C$ is \emph{finitely pairwise
compositional} if there exists a symmetric polynomial $P \in
\mathbb{R}[u, v]$ of total degree at most two such that the derived cost
satisfies
\[
F(xy) + F\!\left(\tfrac{x}{y}\right) = P(F(x), F(y))
\quad \text{for all } x, y > 0.
\]
\end{definition}

This is the regularity hypothesis (R2) of Section~\ref{sec:conditions},
restated cleanly: the combiner is a symmetric polynomial of total degree
at most two on $\mathbb{R}^2$. The finite pairwise polynomial closure
condition is the structural restriction we name and verify against the
counterexample below.

\subsection{The direct theorem}

\begin{theorem}[RCL as the finite pairwise algebra of positive-ratio comparison]
\label{thm:rcl-direct}
Let $C$ be an operative positive-ratio comparison that is finitely
pairwise compositional. Then there exists $c \in \mathbb{R}$ such that
the derived cost $F(r) = C(r, 1)$ satisfies the Recognition Composition
Law
\[
F(xy) + F\!\left(\tfrac{x}{y}\right) = 2F(x) + 2F(y) + c\,F(x) F(y)
\quad \text{for all } x, y > 0. \tag{$\mathrm{RCL}$}
\]
\end{theorem}

\begin{proof}
By Identity, $F(1) = C(1, 1) = 0$. By Scale invariance, for any
$x, y > 0$,
\[
C(x, y) = C(x/y, 1) = F(x/y).
\]
By Non-contradiction together with Scale invariance,
\[
F(x) = C(x, 1) = C(1, x) = C(x^{-1} \cdot 1, x^{-1} \cdot x) = C(x^{-1}, 1) = F(x^{-1}),
\]
so $F$ is reciprocal-symmetric.

By finite pairwise polynomial closure, there exists a symmetric polynomial
$P \in \mathbb{R}[u, v]$ of total degree at most two with
\[
F(xy) + F(x/y) = P(F(x), F(y)).
\]
Set $y = 1$. Since $F(1) = 0$,
\[
F(x) + F(x) = P(F(x), F(1)) = P(F(x), 0),
\]
so $P(F(x), 0) = 2 F(x)$ for all $x > 0$. Because $F$ is continuous and
nontrivial on the connected domain $\mathbb{R}_{>0}$, the image
$\mathrm{Range}(F)$ contains a nondegenerate interval $I \subseteq
\mathbb{R}$. The polynomial identity
\[
P(u, 0) = 2u
\]
therefore holds on $I$, and by polynomial unique continuation it holds
for all $u \in \mathbb{R}$. By symmetry of $P$,
\[
P(0, v) = 2v \qquad \text{for all } v \in \mathbb{R}.
\]

Write the most general symmetric polynomial of total degree at most two:
\[
P(u, v) = a + b\,(u + v) + d \, uv + e\,(u^2 + v^2),
\qquad a, b, d, e \in \mathbb{R}.
\]
Using $P(u, 0) = 2u$,
\[
a + b\,u + e\,u^2 = 2u
\quad \text{for all } u \in \mathbb{R},
\]
which forces $a = 0$, $b = 2$, and $e = 0$. Hence
\[
P(u, v) = 2u + 2v + d \, uv.
\]
Renaming $d = c$, the conclusion follows:
\[
F(xy) + F(x/y) = 2 F(x) + 2 F(y) + c\,F(x) F(y).
\qedhere
\]
\end{proof}

\begin{remark}[Comparison with the d'Alembert route]
\label{rem:direct-vs-dAlembert}
The proof of Theorem~\ref{thm:rcl-direct} is shorter and more
self-contained than the d'Alembert route used in
Theorem~\ref{thm:main}. The two routes give the same conclusion. The
d'Alembert route routes through a published rigidity
result~\cite{WashburnZlatanovicAllahyarov2026} that handles a wider
class of related functional equations and locates the present case as
one instance. The direct route uses only polynomial unique continuation
on a nondegenerate interval. We retain both because the d'Alembert
route gives broader context and the direct route gives self-containment.
\end{remark}

\subsection{Necessity of finite pairwise polynomial closure}

The hypothesis of finite pairwise polynomial closure
(Definition~\ref{def:finite-pairwise-closure}) is essential. We
demonstrate this by exhibiting a comparison operator that satisfies the
operative positive-ratio conditions and admits a continuous combiner,
but whose combiner is not polynomial of degree two and produces a cost
function that does not satisfy the Recognition Composition Law.

\begin{proposition}[Necessity of finite pairwise polynomial closure]
\label{prop:finite-closure-necessary}
The conclusion of Theorem~\ref{thm:rcl-direct} fails if finite pairwise
polynomial closure is replaced by arbitrary continuous compositionality.
\end{proposition}

\begin{proof}
Define the comparison operator
\[
C(x, y) := \bigl(\ln(x/y)\bigr)^4,
\qquad F(x) := C(x, 1) = (\ln x)^4.
\]
The operator is continuous on $\mathbb{R}_{>0}^2$, is total on the
positive quadrant, satisfies Identity ($C(x, x) = 0$), is symmetric
($C(x, y) = (\ln(x/y))^4 = (\ln(y/x))^4 = C(y, x)$), and is scale
invariant ($C(\lambda x, \lambda y) = (\ln(\lambda x / \lambda y))^4 =
(\ln(x/y))^4 = C(x, y)$). Hence $C$ is an operative positive-ratio
comparison in the sense of Definition~\ref{def:operative-comparison}.

Set $t = \ln x$ and $u = \ln y$. Then
\[
F(xy) + F(x/y) = (t + u)^4 + (t - u)^4 = 2 t^4 + 12 t^2 u^2 + 2 u^4.
\]
Since $F(x) = t^4$ and $F(y) = u^4$, we have $t^2 u^2 = \sqrt{F(x) F(y)}$,
so
\[
F(xy) + F(x/y) = 2 F(x) + 2 F(y) + 12 \sqrt{F(x) F(y)}.
\]
The combining rule
\[
\Phi(a, b) = 2 a + 2 b + 12 \sqrt{a b}
\]
is a continuous symmetric function on the range
$\mathrm{Range}(F) \times \mathrm{Range}(F) = [0, \infty)^2$, and the
operator $C$ satisfies the d'Alembert form $F(xy) + F(x/y) = \Phi(F(x), F(y))$.

If the Recognition Composition Law held for $F$, there would exist a
constant $c \in \mathbb{R}$ such that
\[
12 \sqrt{a b} = c \, a b
\quad \text{for all } a, b > 0.
\]
Setting $a = b = 1$ gives $c = 12$. Setting $a = b = 4$ then gives
$c = 12 \cdot 2 / 16 = 3 / 2$, contradicting $c = 12$. Hence no constant
$c$ works, and the Recognition Composition Law fails for $F$.

The combining rule $\Phi(a, b) = 2 a + 2 b + 12 \sqrt{a b}$ is not a
polynomial of total degree at most two on $\mathbb{R}^2$ (the $\sqrt{a b}$
term is not even smooth at the origin). The example therefore shows
that an operative positive-ratio comparison can admit a continuous
combiner without the combiner being polynomial of degree at most two,
and that in such a case the Recognition Composition Law is not
satisfied. The finite pairwise polynomial closure hypothesis is
therefore essential, not redundant; it is the precise structural
condition that distinguishes the Recognition Composition Law family
from broader continuous-combiner alternatives.
\end{proof}

\begin{remark}[The status of finite pairwise polynomial closure]
\label{rem:r2-status-update}
Proposition~\ref{prop:finite-closure-necessary} reframes the role of
the polynomial-degree-two restriction (R2). It is not a regularity
hypothesis carried alongside the four classical conditions for technical
convenience and hopefully removable. It is the structural content that
distinguishes the Recognition Composition Law family from broader
algebras of continuous positive-ratio comparison such as the
$(\ln(x/y))^4$ family. Removing (R2) does not produce a more general
form of RCL; it produces a strictly larger class of compositional
algebras that includes structures (such as the quartic-log example) not
governed by the Recognition Composition Law. The polynomial restriction
is therefore part of the structural meaning of \emph{finite pairwise
algebra} in the present setting. Section~\ref{subsec:conjecture} records
why two attempted weakenings, closure under iteration and analyticity at
the origin, both fail to recover the polynomial restriction.
\end{remark}

\section{Calibrated Costs on the Bilinear and Additive Branches}\label{sec:cost}

Theorem~\ref{thm:main} fixes the polynomial form of the route-independence
combiner up to one real parameter $c$. To pin down $c$, we add the
non-negativity, convexity, and curvature conditions standard in the
characterization of the canonical reciprocal cost. Both the bilinear
($c \neq 0$) and additive ($c = 0$) branches survive these additional
conditions, each with its own canonical calibrated form.

\begin{definition}[Convex non-negative comparison operator]
\label{def:convex-nonnegative}
A comparison operator $C$ is \emph{convex non-negative} if its derived cost
$F$ is non-negative and convex on $\mathbb{R}_{>0}$.
\end{definition}

\begin{remark}
Non-negativity is the natural condition for a cost interpretation: a
comparison's cost is a real-valued penalty for deviation from equilibrium,
not a signed quantity. Convexity is the natural regularity demand of
penalty functions: the cost of intermediate ratios is bounded by an average
of the costs at the endpoints. Both conditions are standard in the
information-geometric literature on ratio-based costs and are part of the
hypotheses of the Canonical Reciprocal Cost Uniqueness
Theorem~\cite{WashburnZlatanovic2026}.
\end{remark}

\begin{definition}[Log-curvature calibration]
\label{def:calibration}
The \emph{log-curvature} of a function $F : \mathbb{R}_{>0} \to \mathbb{R}$
at the multiplicative identity is
\[
\kappa(F) := \lim_{t \to 0} \frac{2F(e^t)}{t^2},
\]
provided the limit exists. We say $F$ satisfies the \emph{unit log-curvature
calibration} if $\kappa(F) = 1$.
\end{definition}

\begin{remark}
The log-curvature condition fixes the second-order behaviour of $F$ at the
identity in logarithmic coordinates. It is the minimal regularity that
isolates a single member of the bilinear family
$P(u, v) = 2u + 2v + c\,uv$. See
\cite[Section~5]{WashburnZlatanovicAllahyarov2026} and
\cite[Section~2]{WashburnZlatanovic2026} for the role of this calibration in
canonical-coefficient selection.
\end{remark}

\begin{theorem}[Calibrated cost on the bilinear branch: the $\alpha$-family
and its canonical representative]
\label{thm:cost-corollary}
Let $C$ be a comparison operator satisfying the encoded conditions
(Definition~\ref{def:laws-of-logic}), with derived cost $F$ convex,
non-negative, and satisfying the unit log-curvature calibration
$\kappa(F) = 1$. Assume in addition that $F$ lies on the bilinear
branch of Theorem~\ref{thm:main}: the parameter $c$ given by that
theorem is non-zero. Then there exists $\alpha \geq 1$ such that
\[
F(x) = F_\alpha(x)
     := \frac{1}{\alpha^2}\!\left(\cosh(\alpha \ln x) - 1\right)
\quad \text{for all } x > 0,
\qquad \text{with } c = 2\alpha^2.
\]
After fixing the multiplicative-coordinate normalization $\alpha = 1$,
the canonical representative of the $\alpha$-family is the canonical
reciprocal cost
\[
J(x) := F_1(x) = \tfrac{1}{2}\!\left(x + x^{-1}\right) - 1,
\qquad c = 2.
\]
\end{theorem}

\begin{proof}
By Theorem~\ref{thm:main}, $F$ satisfies the bilinear composition law
$F(xy) + F(x/y) = 2F(x) + 2F(y) + c\,F(x) F(y)$ for some $c \in \mathbb{R}$,
and by hypothesis $c \neq 0$. By Theorem~\ref{thm:translation}, $F$ is
continuous and nonconstant with $F(1) = 0$. By
Lemma~\ref{lem:reciprocity}, $F(x) = F(x^{-1})$. By hypothesis $F$ is
convex, non-negative, and satisfies $\kappa(F) = 1$.

By
\cite[Corollary~8]{WashburnZlatanovicAllahyarov2026}, convexity and
non-negativity in the bilinear branch $c \neq 0$ force the hyperbolic
form $F(e^t) = (2/c)(\cosh(\alpha t) - 1)$ for some $\alpha \geq 1$,
with $c > 0$. By \cite[Theorem~9]{WashburnZlatanovicAllahyarov2026},
the unit log-curvature calibration $\kappa(F) = 1$ implies
$c = 2\alpha^2$. Substituting,
\[
F(e^t) = \frac{2}{c}\,(\cosh(\alpha t) - 1)
       = \frac{1}{\alpha^2}\,(\cosh(\alpha t) - 1),
\]
which yields the $\alpha$-family $F = F_\alpha$ as stated. After the
multiplicative-coordinate normalization $\alpha = 1$
(\cite[Corollary~10]{WashburnZlatanovicAllahyarov2026}), $c = 2$ and
\[
F(x) = \cosh(\ln x) - 1
     = \tfrac{1}{2}(x + x^{-1}) - 1
     = J(x).
\]
Equivalently, the canonical representative ($\alpha = 1$) coincides
with the unique continuous reciprocal-symmetric solution of the
canonical $c = 2$ composition law on $\mathbb{R}_{>0}$ with $F(1) = 0$
and $\kappa(F) = 1$, by the Canonical Reciprocal Cost Uniqueness
Theorem~\cite[Theorem~1]{WashburnZlatanovic2026}.
\end{proof}

\begin{remark}[Self-application on the positive-real domain]
\label{rem:self-application}
The canonical reciprocal cost $J$, and the $\alpha$-family more
generally, is a function $\mathbb{R}_{>0} \to [0, \infty)$. On inputs
$x \neq 1$ the output is strictly positive, and any positive output
lies in the domain $\mathbb{R}_{>0}$ on which $J$ itself is defined.
Consequently, for any $x, y > 0$ with $J(x), J(y) > 0$, the ratio
$J(x)/J(y)$ is a positive real and $J(J(x)/J(y))$ is a well-defined
element of $[0, \infty)$. The functional form is closed under its own
composition on non-trivial inputs. The four encoded conditions
(L1)--(L4) on a comparison operator apply verbatim at this meta-level:
the cost of comparing two costs satisfies Identity, single-valuedness
under reordering, totality on the open positive quadrant, and
composition consistency by the same functional form as at the base
level. The self-application is not an additional hypothesis; it is a
direct consequence of $\mathbb{R}_{>0}$ being both the domain and the
positive part of the codomain. On the restricted domain, the
derived-cost operator is its own comparison operator: the structural
conditions that characterise it at any level characterise it at every
level.
\end{remark}

\begin{remark}[Why $J$ is not directly forced]\label{rem:why-J-not-forced}
The unit log-curvature calibration alone does not force $c = 2$ on
the bilinear branch. It forces $c = 2\alpha^2$ for some $\alpha \geq 1$.
The canonical reciprocal cost $J$ is the $\alpha = 1$ member of the
family: the multiplicative-coordinate scale at which the cost rises
quadratically in $\ln x$ at unit rate near the identity. Different
$\alpha$ correspond to rescaling $x \mapsto x^\alpha$; the family
$F_\alpha$ is the orbit of $J$ under this rescaling. The choice
$\alpha = 1$ is a coordinate normalisation, not a consequence of the
four classical laws plus the regularity hypotheses. Treating it as one
would conflate the classical laws with a coordinate convention. The
theorem above states the $\alpha$-family explicitly and identifies $J$
only as the normalised representative.
\end{remark}

\begin{remark}[The additive branch and the choice between branches]
\label{rem:additive-branch}
If instead $c = 0$ in the conclusion of Theorem~\ref{thm:main}, the
composition law reduces to $F(xy) + F(x/y) = 2F(x) + 2F(y)$. Continuous
solutions with $F(1) = 0$ have the form $F(x) = k(\ln x)^2$ for some
$k \in \mathbb{R}$
\cite[Theorem~5(ii)]{WashburnZlatanovicAllahyarov2026}. The unit
log-curvature calibration $\kappa(F) = 1$ forces $k = 1/2$, giving
\[
F(x) = \tfrac{1}{2}(\ln x)^2,
\]
which is convex, non-negative, and reciprocal-symmetric. This is the
parallel canonical form on the additive branch, distinct from $J$ but also
unique within its branch under the same calibration. The two canonical
forms correspond to the two branches of the d'Alembert classification on
$\mathbb{R}_{>0}$ (see
\cite[Remark~9]{WashburnZlatanovicAllahyarov2026} and
\cite[Section~3]{WashburnRahnamai2026}). The four Aristotelian conditions,
together with scale invariance, non-triviality, convexity, non-negativity,
and the calibration $\kappa(F) = 1$, do not by themselves distinguish the
two branches: each produces a unique canonical form. What further
structural condition selects the bilinear branch $c \neq 0$ over the
additive branch $c = 0$ is an open question; we return to it in
Section~\ref{subsec:additive-branch}.
\end{remark}

\begin{remark}[The Reciprocal Convex Cost characterization]
A complementary characterization of $J$ in the optimization-theoretic
setting was established by~\cite{WashburnRahnamai2026}, where ratio-induced
cost functions of the form $c(s, o) = J(\iota_S(s)/\iota_O(o))$ are shown,
under inversion symmetry, strict convexity, coercivity, normalization at
the identity, and a multiplicative d'Alembert identity, to admit the
hyperbolic family $J(x) = \tfrac{1}{2}(x^a + x^{-a}) - 1$ as the unique
admissible mismatch penalty, with $a > 0$. The present cost theorem and
the optimization-theoretic characterization land in the same place
because they invoke compatible regularity hypotheses (continuity, convexity,
multiplicative d'Alembert composition with appropriate calibration). We
therefore characterize this not as three independent derivations
converging on the same object, but as three readings of the same
hypothesis class; we discuss the encoding-dependence of this convergence
in Section~\ref{sec:wigner}.
\end{remark}

\section{The Operative Domain: Reality, Logic, and the Recognition Composition Law}\label{sec:operative}

The translation theorem (Section~\ref{sec:translation}), the main
theorem (Section~\ref{sec:main}), and the calibrated cost theorem
(Section~\ref{sec:cost}) together establish the rigidity of the
recognition composition law and its calibrated solutions on a
continuous comparison operator satisfying the encoded conditions. The
canonicality theorem (Section~\ref{sec:canonicality}) shows that the
encoding is the unique reading of the four classical Aristotelian
conditions under the magnitude-of-mismatch interpretation. We now
identify the scope on which these results land as a single structural
identification, and we show that the scope is the operative domain of
physics and continuous mathematics on positive quantities.

\subsection{Reality structures}

\begin{definition}[Reality structure]\label{def:reality-structure}
A \emph{reality structure} is a triple $(\mathcal{S}, C, \mathcal{R})$
where $\mathcal{S}$ is a non-empty set of states,
$C : \mathcal{S} \times \mathcal{S} \to \mathbb{R}$ is a comparison
operator returning real-valued mismatches between states, and
$\mathcal{R}$ is a relation of compositional inference on
$(\mathcal{S}, C)$. The triple is a reality structure if the four
following structural conditions hold.
\begin{enumerate}[label=(R\arabic*'),leftmargin=*]
\item \emph{Definite states.} For every $s \in \mathcal{S}$, the state
$s$ is a single, well-defined element. Equivalently, $\mathcal{S}$ is
a set in the standard sense of mathematics.
\item \emph{Single-valued comparison.} $C(s_1, s_2)$ returns a single
real value for every pair $(s_1, s_2) \in \mathcal{S} \times
\mathcal{S}$, and the value does not depend on the side from which the
comparison is read.
\item \emph{Decidability.} For every pair $(s_1, s_2) \in \mathcal{S}
\times \mathcal{S}$, $C(s_1, s_2)$ is defined: there is no input pair
on which the comparison returns ``undefined'' or some third value
outside $\mathbb{R}$.
\item \emph{Compositional inference.} For composite comparisons
assembled by the relation $\mathcal{R}$, the mismatch of the composite
is determinable from the mismatches of its constituents.
\end{enumerate}
\end{definition}

The four conditions (R1$'$)--(R4$'$) are the structural minimum any
object can satisfy and still be called ``a reality'' in the standard
usage of the word. They are weaker than the conditions of an ordinary
structure in mathematical logic. A Boolean algebra is a reality
structure (with $C$ a discrete-valued mismatch indicator); a metric
space is a reality structure (with $C$ the metric); the standard
ZFC universe is a reality structure (with $C$ any decidable
comparison operator); physical reality, modelled as the structure of
states accessible to measurement, is a reality structure under the
standard reading of measurement as a comparison of one state to
another with a real-valued mismatch.

\begin{theorem}[Reality satisfies the four classical conditions]
\label{thm:reality-satisfies-logic}
Any reality structure $(\mathcal{S}, C, \mathcal{R})$ satisfies the
four classical Aristotelian conditions on its comparison operator: $C$
satisfies Identity, Non-contradiction, Totality, and Composition
consistency.
\end{theorem}

\begin{proof}
By definition unfolding.
\begin{itemize}[leftmargin=*]
\item (R1$'$) implies (L1): a single, well-defined state is the same
as itself; the mismatch of a state with itself is the trivial value of
the codomain, namely $0$. Thus $C(s, s) = 0$ for every
$s \in \mathcal{S}$.
\item (R2$'$) implies (L2): single-valued comparison forbids two
distinct values for the same comparison; the value does not depend on
the side from which the comparison is read, so $C(s_1, s_2) = C(s_2,
s_1)$.
\item (R3$'$) implies (L3): decidability is exactly totality on
$\mathcal{S} \times \mathcal{S}$.
\item (R4$'$) implies (L4): determinability of composite mismatches
from constituents is exactly the existence of a combining rule for
composite costs.
\end{itemize}
The four logical-content conditions hold on any reality structure.
\end{proof}

\subsection{The operative domain}

\begin{definition}[Operative domain]
\label{def:operative-domain}
A reality structure $(\mathcal{S}, C, \mathcal{R})$ is
\emph{operative} if $\mathcal{S} \subseteq \mathbb{R}_{>0}$ (states
are encoded as positive reals), $C$ is jointly continuous on
$\mathcal{S} \times \mathcal{S}$ (regularity hypothesis (R1) holds),
and the combining rule $\Phi$ given by (R4$'$) is polynomial of
degree at most two on $\mathbb{R}^2$ (regularity hypothesis (R2)
holds).

The \emph{operative domain} is the class of operative reality
structures.
\end{definition}

The operative domain is broad. Every physical measurement is a
comparison of two positive quantities (the magnitudes of two states,
or a state and a calibrated reference), with a continuous real-valued
cost capturing the discrepancy. Every continuous likelihood ratio
between two probability distributions is a comparison on the
operative domain. Every continuous distance or divergence between two
states in a continuous state space is a comparison on the operative
domain. The mathematical structure of the natural sciences in their
pairwise-comparison setting lives in the operative domain.

\begin{theorem}[The operative-domain identification]
\label{thm:operative-identification}
On the operative domain, the recognition composition law
$F(xy) + F(x/y) = 2F(x) + 2F(y) + cF(x)F(y)$ is the law of logic.
More precisely:
\begin{enumerate}[label=(\alph*),leftmargin=*]
\item Any operative reality structure has its derived cost $F$
satisfying the recognition composition law for some
$c \in \mathbb{R}$.
\item Under convexity, non-negativity, and the unit log-curvature
calibration, the bilinear branch ($c \neq 0$) admits the calibrated
representative $J(x) = \tfrac{1}{2}(x + x^{-1}) - 1$, and the additive
branch ($c = 0$) admits the calibrated representative
$\tfrac{1}{2}(\ln x)^2$.
\item The encoding from the structural conditions (R1$'$)--(R4$'$) of
a reality structure to the four logical-content conditions
(L1)--(L4) is canonical under the magnitude-of-mismatch
interpretation of the comparison operator.
\end{enumerate}
\end{theorem}

\begin{proof}
(a) Any operative reality structure satisfies (R1$'$)--(R4$'$) by
Definition~\ref{def:reality-structure}; thus by
Theorem~\ref{thm:reality-satisfies-logic} the comparison operator
satisfies (L1)--(L4); the operative-domain definition supplies the
regularity hypotheses (R1) and (R2); the bridge condition (B1) is
imposed structurally on $\mathcal{S} \subseteq \mathbb{R}_{>0}$ in
the operative-domain definition (states differing by a multiplicative
constant are identified at the comparison level under the standard
cost reading); thus by Theorem~\ref{thm:main} the derived cost
satisfies the recognition composition law for some
$c \in \mathbb{R}$.

(b) is Theorem~\ref{thm:cost-corollary} and
Remark~\ref{rem:additive-branch}.

(c) is Theorem~\ref{thm:canonicality}.
\end{proof}

\subsection{The chain on the operative domain}

The three theorems together identify, on the operative domain, three
distinct objects as a single mathematical entity. We state the
identification.

\begin{corollary}[Recognition composition law $=$ law of logic $=$ structural form of reality, on the operative domain]
\label{cor:rcl-logic-reality}
On the operative domain:
\begin{enumerate}[label=(\arabic*),leftmargin=*]
\item The structural form of a reality
(Definition~\ref{def:reality-structure}, restricted to the operative
domain) is the four classical Aristotelian conditions on the
comparison operator (Theorem~\ref{thm:reality-satisfies-logic}).
\item The unique mathematical object compatible with the four
classical Aristotelian conditions, under the canonical encoding, is
the recognition composition law and its calibrated solutions
(Theorems~\ref{thm:canonicality}, \ref{thm:translation},
\ref{thm:main}, \ref{thm:cost-corollary}).
\item Therefore the structural form of a reality on the operative
domain is the recognition composition law: any reality structure on
the operative domain has its comparison operator's derived cost
satisfying RCL, with the calibrated representatives $J$ on the
bilinear branch and $\tfrac{1}{2}(\ln x)^2$ on the additive branch.
\end{enumerate}
\end{corollary}

The chain ``recognition composition law $=$ law of logic $=$
structural form of reality'' is therefore a theorem on the operative
domain. The identification is not a structural correspondence under
one encoding among many; it is the unique consequence of (i) the
canonical reading of the classical Aristotelian conditions on a
continuous comparison operator (Theorem~\ref{thm:canonicality}),
(ii) the structural definition of a reality as a single-valued,
decidable, compositionally consistent comparison structure
(Definition~\ref{def:reality-structure}), and (iii) the rigidity of
the d'Alembert combiner under the regularity hypotheses (R1) and (R2)
(Theorem~\ref{thm:main}).

\subsection{Scope and remaining open questions}

The corollary establishes the identification on the operative domain.
Three classes of questions remain open.

\paragraph{The discrete propositional case.} A discrete reality
structure (one in which $\mathcal{S}$ is finite or countable, and $C$
is a discrete-valued comparison) does not directly fall under the
operative domain. However, a standard route via Stone's representation
theorem~\cite{Stone1936} embeds any Boolean algebra into the algebra
of clopen subsets of a totally disconnected compact Hausdorff space,
and measure-theoretic completion lifts the discrete-valued comparison
to a continuous-valued likelihood ratio between probability
distributions on the resulting space. Under this lifting, the discrete
propositional case is a continuous limit of the operative domain.
Whether the recognition composition law applies in its native form to
the discrete case, or whether it admits a discrete cousin (a
polynomial of degree at most two on a discrete subset of
$\mathbb{R}$), is the natural extension of the present paper. We do
not settle it here.

\paragraph{Modal, intuitionistic, and higher-order logics.} The
encoding (L1)--(L4) presented in this paper is for classical
Aristotelian logic on a continuous comparison operator. Modal logics,
intuitionistic logic, and higher-order logics may admit different
encodings of their structural conditions, with potentially different
rigidity results. We do not address them here.

\paragraph{The metaphysical reading of ``reality structure.''}
Definition~\ref{def:reality-structure} of a reality structure is
structural: it abstracts the four properties (definite states,
single-valued comparison, decidability, compositional inference) as
the structural minimum any object can satisfy and still be called ``a
reality'' in the standard usage. The metaphysical reading of
``reality'' --- what counts as a reality in the world, in some
metaphysical sense beyond the structural minimum --- is a separate
question that we leave outside the scope of the present paper. The
present result is structural: on any structure satisfying the four
properties, in the operative-domain regularity setting, the
recognition composition law is the law of the comparison operator.

\section{Interpretation: The Encoding, Wigner's Question, and the Historical Separation}\label{sec:wigner}

This section is interpretive. The mathematical content of the paper is
the Translation Theorem and the rigidity invocations of
Sections~\ref{sec:translation}--\ref{sec:cost}. What follows is the
discussion of what those theorems do and do not establish about the
relationship between logical principles and mathematical objects. We
keep the philosophical claims modest and explicitly tied to the
encoding.

\subsection{What the prior formalizations did}

The formalizations of logic listed in the introduction, from
Aristotle's syllogistic to univalent foundations, are formal systems
within which valid reasoning is described. Each provides a calculus, a
notation, an algebra, or a semantic framework. Each makes precise the
behaviour of logic under specified operational interpretations.

These articulations are not displaced by the present result. Each
remains correct as an account of how the laws of logic behave under the
chosen interpretation. What the present paper adds, on a restricted
domain, is a structural correspondence: under the encoding
(L1)--(L4) of Section~\ref{sec:conditions}, augmented by the named
regularity hypotheses (R1) and (R2), the four classical conditions and
the Recognition Composition Law are co-extensive. They constrain the
same comparison operators.

The Feynman ``name versus thing'' analogy~\cite{Feynman1964} is useful
here, but cuts both ways. The prior formalizations give the laws of
logic increasingly precise names, in increasingly expressive notations,
without exhibiting the laws as a forced functional-equation structure
on continuous comparisons. The present result, on its restricted
domain, identifies such a structure: a single forced combiner family
with which the encoded conditions are co-extensive, with two
calibrated branch representatives.
Whether this identification \emph{is} the bird's anatomy, or whether it
is itself another name (the encoding being a labeling of mathematical
conditions with Aristotelian-style vocabulary), is a substantive
question that the encoding-defenses of
Section~\ref{subsec:encoding-defenses} address but do not foreclose. We
discuss this directly below.

\subsection{Wigner on the unreasonable effectiveness of mathematics}

In 1960, Eugene Wigner posed the puzzle of ``the unreasonable
effectiveness of mathematics in the natural sciences''~\cite{Wigner1960}.
Mathematics, conducted historically in pursuit of internal coherence,
gives quantitative descriptions of physical phenomena that exceed the
precision of the experiments testing them. The puzzle is why such an
agreement should hold at all.

Standard responses reframe the puzzle without dissolving it.
\emph{Mathematics is a language of patterns}: this restates the
agreement; it does not explain it. \emph{The universe is lawful and
lawfulness is what mathematics describes}: this assumes that lawfulness
takes the form of the mathematical language we use, without independent
reason. \emph{We evolved to recognise the patterns the universe
presents}: this addresses our ability to recognise the agreement, not
why the agreement is there.

\subsection{What the present result does and does not do for Wigner's
question}\label{subsec:wigner-honest}

We address the question carefully. The result of the present paper is
that, under the proposed encoding, the same continuous functional
equation governs (a) the rigidity of polynomial composition combiners
for continuous ratio costs (the d'Alembert Inevitability Theorem
\cite{WashburnZlatanovicAllahyarov2026}) and (b) the encoded conditions
(L1)--(L4) plus regularity hypotheses (R1), (R2) on a comparison
operator. Both characterizations land on the Recognition Composition
Law family, and on the canonical reciprocal cost $J$ after branch
selection and coordinate normalization.

The skeptic's objection is direct. Are these two characterizations
independent, or is the convergence built into the encoding? Expanded,
the encoded conditions are: $C(x, x) = 0$ (which gives $F(1) = 0$),
$C(x, y) = C(y, x)$ (which gives reciprocity given scale invariance),
joint continuity of $C$ (which gives continuity of $F$), the polynomial
composition law on $F$, and non-triviality of $F$. These are the
hypotheses of the d'Alembert Inevitability Theorem with
Aristotelian-style labels attached. The convergence is between two
readings of the same hypotheses, not between two derivations from
disjoint first principles.

The skeptic is right. The paper does not establish that two
mathematically independent derivations converge on $J$. What it
establishes is narrower: one natural Aristotelian-style reading of the
hypotheses required by the d'Alembert Inevitability Theorem on a
continuous comparison operator exists, and can be written down
precisely. That is structural information about the domain. It does
not dissolve Wigner's puzzle. It gives one case where the same
hypotheses admit two readings, one of which is classical-logical.

Whether relabelings of this kind exist on other domains (discrete
logic, modal logic, propositional calculus on Boolean algebras) is
the natural extension of the project. The present paper identifies the
phenomenon on one domain and does not claim to answer Wigner's
question.

\subsection{The historical separation}

Philosophers and mathematical logicians developed the laws of logic
for analysing valid reasoning. Physicists and information theorists
developed ratio-based cost functions for modelling penalties for
deviation from equilibrium. The two traditions used different
vocabularies. The present result exhibits, on a restricted domain, a
single forced combiner family with which both vocabularies engage.

This is not a foundational reduction. It does not show classical logic
is reducible to cost-physics or cost-physics to classical logic. It
shows that on a continuous comparison operator on positive ratios,
both traditions describe the same mathematical conditions, and those
conditions admit a unique combiner family with two unique calibrated
solutions, one per branch. The identification is conditional on the
encoding of Section~\ref{sec:conditions}; alternative encodings might
land elsewhere or fail to converge on a single object.

Russell and Whitehead's
\emph{Principia Mathematica}~\cite{RussellWhitehead1913} sought a
foundational reduction; G\"odel's incompleteness
theorems~\cite{Goedel1931} circumscribed any such reduction. The
present paper is not such a reduction. It is one instance, at the
level of a single continuous functional equation, of the connection
between the mathematical and the logical on a specific domain.

\subsection{Implications for long-standing open questions}\label{subsec:implications}

Conditional on the encoding being judged adequate, the rigidity result
has bearing on three long-standing open questions in the philosophy of
mathematics and logic. Each implication below is discussion, not
theorem. Each weakens if the encoding is judged inadequate.

\paragraph{The mathematical-Platonism debate.}
The philosophy of mathematics has carried an unsettled debate for
roughly twenty-four centuries about whether mathematical objects are
discovered (Platonism, Pythagoreanism, mathematical realism) or
invented (formalism, conventionalism, fictionalism). The present
result is relevant in one narrow sense: on the continuous-ratio
comparison domain, the mathematical object characterized by the four
encoded conditions is the same object characterized by
ratio-based-cost rigidity considerations independent of formalisation.
The mathematician's freedom to choose a different object on this
domain is constrained. Under the named hypotheses, no other continuous
object is admissible. Whether this counts as a partial vindication of
mathematical realism or as a structural observation about which
objects are admissible depends on the philosophical framework brought
to it. The rigidity result furnishes one example, on a restricted
domain, of a mathematical object whose form is forced rather than
chosen.

\paragraph{Mathematical intuition and physical reality.}
Wigner's puzzle, addressed in Section~\ref{subsec:wigner-honest}, is a
particular instance of a broader question: why does the mathematician's
intuition about which mathematical objects are natural agree, with such
accuracy, with the experimental physicist's findings about which
objects describe physical reality? The present result provides one
structural account of the agreement, on the continuous-ratio comparison
domain. The object physical-cost rigidity selects is the same object the
encoded Aristotelian conditions admit on a continuous comparison
operator. The mathematician's intuition that $J$ is a natural object on
positive ratios, and the experimental physicist's identification of $J$
as the object physical-cost rigidity selects, are not two independent
confirmations of the same conclusion: they are two readings of the same
hypothesis class. On this domain, mathematical intuition tracks physical
reality because both are reading the same set of structural conditions.
The relationship is structural, not metaphysical. It does not establish
that mathematical intuition tracks physical reality on every domain; it
shows that one mechanism by which the agreement could obtain (the same
hypotheses governing both) does in fact obtain on this restricted
domain.

\paragraph{The foundations of mathematics.}
The foundational programmes of the late nineteenth and early twentieth
centuries, namely logicism (Frege, Russell--Whitehead), formalism
(Hilbert), and intuitionism (Brouwer, Heyting), each sought a unique
foundation for mathematics in either logic, formal axioms, or
constructive proof. G\"odel's incompleteness
theorems~\cite{Goedel1931} circumscribed any such reduction.
Tarski~\cite{Tarski1933}, Gentzen~\cite{Gentzen1935}, the Curry--Howard
correspondence~\cite{Curry1934, Howard1969}, Lawvere's
categorical logic~\cite{Lawvere1964}, Girard's linear
logic~\cite{Girard1987}, and the univalent foundations
program~\cite{HoTT2013} have each provided foundational frameworks of
different syntactic and semantic content. Each is correct as far as it
reaches; none is universally accepted as the foundation. The present
result does not settle the foundations question. It observes, on the
continuous-ratio comparison domain, that all the foundational
frameworks listed agree, by their treatment of continuous comparisons
of positive ratios, on a single forced combiner family with the same
two calibrated branch representatives: $J$ on the bilinear branch,
$\tfrac{1}{2}(\ln x)^2$ on the additive branch. The
frameworks differ in syntactic conventions, in extensions beyond the
continuous-ratio domain, and in the choices of which axioms are
fundamental; on the continuous-ratio domain itself they agree on the
object characterized by the four Aristotelian conditions plus the
regularity hypotheses. This is structural information about the
restricted domain. It is consistent with multiple foundational
frameworks being adequate as long as they agree on the restricted-domain
rigidity. It does not select a preferred foundation, and it does not
violate G\"odel's incompleteness, since the result is conditional on
hypotheses (the encoding plus regularity) that are themselves stated in
classical mathematics with its incompleteness intact.

\paragraph{Physical laws as conventions versus features.} The
philosophy of physics has carried, for roughly three centuries, a
debate over whether physical laws are human conventions adopted for
descriptive convenience or features of the world discovered rather
than stipulated. The conventionalist tradition traces through
Poincar\'e and receives twentieth-century articulation in Reichenbach;
the realist tradition carries into contemporary mathematical physics.
The present result is one data point on the restricted domain where
physical laws admit expression as pairwise comparisons on positive
ratios. On that domain, under the encoding, the combiner family
$P(u, v) = 2u + 2v + c\,uv$ is not chosen but forced: no other
combiner family is admissible. This is a structural fact about the
form of the admissible laws on the restricted domain, not a convention
for describing them. The result neither settles the debate nor claims
to address its broader scope. It gives one case where, on a specific
domain, the conventionalist option is mathematically eliminated.

\paragraph{Comparison operators across substrates.} A comparison
operator is an abstract structure whose instantiations can be physical
(a measurement apparatus comparing two quantities), computational (an
algorithm returning a distance between two inputs), or cognitive (a
judgment of similarity between two perceptions). The rigidity result
on the continuous positive-ratio domain applies to any such
instantiation under the encoding and regularity hypotheses stated. We
do not extend this into a metaphysical claim about the relationship
between physical substrate and cognitive substrate. The observation is
structural: on the restricted domain, the conditions governing
comparison admit a unique functional form regardless of the substrate
in which the comparison is realised. To the extent that cognitive
operations can be modelled as continuous positive-ratio comparisons,
an active question in cognitive science and information geometry, the
same canonical calibrated forms govern both physical and cognitive
realisations. Whether this structural agreement bears on the broader
mind-matter question, which has been carried by philosophical
literature from Descartes onward, is beyond the scope of the present
paper.

\paragraph{Boundary behaviour and the existence-consistency question.}
On the continuous-ratio comparison domain, the canonical
bilinear-branch representative $J$ satisfies $J(x) \to \infty$ as
$x \to 0^+$ and as $x \to \infty$, with $J(1) = 0$ as the global
minimum; the additive-branch representative $\tfrac{1}{2}(\ln x)^2$
has the same boundary divergence. The divergence is the structural
reason that limiting configurations are unreachable by cost
minimisation on the positive-ratio domain.

The philosophical question of why there is something rather than
nothing, posed by Leibniz, has at times been addressed by identifying
existence with logical consistency: what is possible is what is
consistent, and the consistent configurations are those that satisfy
the appropriate structural conditions. The present result bears on
this question narrowly. Under the encoding on continuous positive-ratio
comparisons, the boundary points $x = 0$ and $x = \infty$ are
unreachable by finite-cost configurations on either branch; totality
(L3) on the open quadrant $\mathbb{R}_{>0} \times \mathbb{R}_{>0}$
together with the boundary divergence of the calibrated solutions
produces this directly. We mention this to indicate a structural
connection on the restricted domain between the boundary behaviour of
the calibrated cost and the classical identification of existence with
consistency. We do not extend the observation into a derivation of
any metaphysical principle, and we leave the broader philosophical
literature on existence and consistency outside the scope of the
present paper.

\paragraph{Self-application of the cost operator.} A further
implication, noted as Remark~\ref{rem:self-application} in the cost
section, is that on the positive-real domain the derived cost operator
is closed under its own application. The non-zero outputs of $J$ lie
in the domain on which $J$ itself is defined, and the four encoded
conditions (L1)--(L4) apply verbatim at the level of costs-of-costs.
This supplies, on the restricted domain, a structural answer to a
version of Wigner's puzzle at the level of the operator itself: the
cost function does not require an external structure to govern its
outputs, because the outputs fall within its own input domain. The
conditions that characterise it at level zero characterise it at every
higher level. We note this as a structural fact about the operator,
not as a derivation of any broader foundational claim.

\paragraph{A scope reminder.} All of the above is conditional on the
encoding (L1)--(L4) being a faithful operationalization of the four
classical conditions on a continuous comparison operator on positive
ratios. Theorem~\ref{thm:canonicality} establishes that the encoding
is canonical under the magnitude-of-mismatch interpretation of the
comparison operator. Section~\ref{subsec:encoding-defenses} defends
the operator-level conditions; Section~\ref{sec:canonicality} then
shows the encoding is the unique reading of the four classical laws
under the standard cost-based interpretation;
Section~\ref{subsec:wigner-honest} addresses the residual concern
that the canonicality holds under one named interpretation rather
than as a metaphysical universal.

\section{Scope and Open Questions}\label{sec:scope}

The present paper establishes a rigidity theorem on a specific structural
domain. We collect here the explicit limits of the result, the natural
extensions, and the remaining open questions.

\subsection{The polynomial regularity assumption (R2)}

The d'Alembert Inevitability Theorem on which Theorem~\ref{thm:main} is
built~\cite[Theorem~3]{WashburnZlatanovicAllahyarov2026} is stated for
\emph{polynomial} combiners $P$ of total degree at most two. The
polynomiality is a regularity assumption. It is not a direct consequence
of any of the four Aristotelian conditions individually; we have
therefore named it explicitly as the regularity hypothesis (R2)
(Definition~\ref{def:route-independence}) and carried it as a separate
ingredient in the main theorem.

The classical theory of d'Alembert's functional equation does not require
polynomiality. The continuous-combiner classification due to
Acz\'el~\cite{Aczel1966}, Kannappan~\cite{Kannappan2009}, and
Stetk{\ae}r~\cite{Stetkaer2013} shows that the only continuous solutions
of $H(t + u) + H(t - u) = 2H(t) H(u)$ on $\mathbb{R}$ with $H(0) = 1$ are
$H(t) = \cosh(\alpha t)$ or $H(t) = \cos(\alpha t)$ for some $\alpha$.
Under suitable additional hypotheses (convexity and non-negativity for the
hyperbolic branch versus the trigonometric branch, branch selection for
$c \neq 0$ versus $c = 0$, calibration for the multiplicative-coordinate
scale), the conclusion of the bilinear branch is the same as in the
polynomial setting.

The generalized result is stated below as a proposition. The same
branch-and-normalization caveats apply: a calibrated continuous
combiner does not isolate $J$ without bilinear-branch selection and
multiplicative-coordinate normalization.

The continuous-combiner extension has been formalised in the canonical
Lean library as the module
\texttt{Foundation.GeneralizedDAlembert}. The module
introduces a continuous-consistency predicate
\texttt{HasContinuousConsistency} (the analogue of
\texttt{HasMultiplicativeConsistency} from \texttt{Inevitability}
without the polynomial restriction), proves a sequence of boundary
identities and structural lemmas, and reduces the bilinear-family
conclusion to a single named smoothness assumption. The fully proved
content of the module is summarised below; the precise sorry-count and
status of each remaining step appears in
Section~\ref{sec:lean-status}.

\begin{itemize}[leftmargin=*]
\item \emph{Boundary identity at $v = 0$.} Under reciprocal symmetry of
$F$ and the normalisation $F(1) = 0$, the combiner satisfies $\Phi(0,
F(y)) = 2 F(y)$ for every $y > 0$. (Lean theorem
\texttt{Phi\_at\_zero\_left}, fully proved.)
\item \emph{Symmetry of $\Phi$ on the range.} Under reciprocal symmetry
of $F$, the combiner is symmetric on
$\operatorname{Range}(F) \times \operatorname{Range}(F)$. (Lean
theorem \texttt{Phi\_symmetric\_on\_range}, fully proved.)
\item \emph{Symmetry of $\Phi$ on the closure.} If
$\operatorname{Range}(F)$ is dense in some interval and $\Phi$ is
continuous, the symmetry on the range extends by continuity to all of
$\mathbb{R}^2$. (Lean theorem \texttt{Phi\_symmetric\_on\_closure},
fully proved, via a closed-set-on-dense-set argument.)
\item \emph{Iteration substitution.} The Acz\'el iteration identity
$H(2t) = \Phi(H(t), H(t))$ holds whenever $H(0) = 0$ and $H$ satisfies
the generalized d'Alembert equation. (Lean theorem \texttt{H\_two\_t},
fully proved.)
\item \emph{Range contains an interval.} A continuous nonconstant $H$
on $\mathbb{R}$ has range containing a nondegenerate interval. (Lean
theorem \texttt{H\_range\_contains\_interval}, fully proved using the
preconnected-image-implies-order-connected route.)
\end{itemize}

The remaining gap is the analytic content of the Acz\'el--Kannappan--Stetk{\ae}r
classification on the combiner side. The Lean module isolates the gap
as the explicit hypothesis package
\texttt{ContinuousCombinerAnalysisInputs}, with three named components:
finite-order mollifier smoothness on the log-coordinate cost, the
second-derivative identity
$2 G''(t) = \psi(G(t)) \cdot G''(0)$, and the $\psi$-affine completion
$\psi(u) = 2 + c \cdot u$ on the image of the log-cost. The package is
deliberately not derived from continuity alone: the same module
formalises the quartic log-cost $C(x, y) = (\ln(x/y))^4$ as an
explicit refutation of the second-derivative identity. When the
package is supplied, the H-side classification (now the discharged
theorem \texttt{aczel\_kannappan\_continuous\_dAlembert}) and the
positive-ratio bilinear assembly close the conclusion mechanically.

\begin{proposition}[Conditional continuous-combiner extension,
branch-aware version]
\label{prop:continuous-combiner-extension}
\emph{Assuming the explicit hypothesis package
\texttt{ContinuousCombinerAnalysisInputs}
(stated in Appendix~\ref{app:lean}, with the quartic log-cost
counterexample showing that the package is not automatic from
continuity alone).}
Let $C$ be a comparison operator on $\mathbb{R}_{>0} \times
\mathbb{R}_{>0}$ satisfying the four logical-content conditions
(L1)--(L4), Continuity (R1), Scale invariance (B1), and Non-triviality
(B2), but with the polynomial regularity (R2) replaced by the weaker
continuous-combiner condition: there exists a continuous symmetric
function $\Phi : \operatorname{Range}(F)^2 \to \mathbb{R}$ for which
\[
F(xy) + F\!\left(\tfrac{x}{y}\right) = \Phi(F(x), F(y))
\quad \text{for all } x, y > 0,
\]
with $\operatorname{Range}(F)$ containing an open interval. Assume
further that $F$ is convex, non-negative, and satisfies the unit
log-curvature calibration $\kappa(F) = 1$. Then exactly one of the
following holds.
\begin{enumerate}[label=(\arabic*), leftmargin=*]
\item \emph{Bilinear branch.} $\Phi(u, v) = 2u + 2v + c\,uv$ on
$\operatorname{Range}(F)^2$ for some $c \neq 0$, $F$ has the form
$F(x) = (1/\alpha^2)(\cosh(\alpha \ln x) - 1)$ for some $\alpha \geq 1$
with $c = 2\alpha^2$, and after fixing the multiplicative-coordinate
normalization $\alpha = 1$ we obtain $c = 2$ and $F = J$.
\item \emph{Additive branch.} $\Phi(u, v) = 2u + 2v$ (the case $c = 0$),
and $F(x) = \tfrac{1}{2}(\ln x)^2$.
\end{enumerate}
\end{proposition}

The proof reduces to the classical Acz\'el--Kannappan--Stetk{\ae}r
classification by passing to logarithmic coordinates. The dichotomy
between the two branches matches Theorem~\ref{thm:cost-corollary} and
Remark~\ref{rem:additive-branch}; the calibration $\kappa(F) = 1$ does
not isolate $J$ without (i) branch selection ($c \neq 0$) and (ii)
coordinate normalization ($\alpha = 1$), both of which are choices
outside the four logical-content conditions.

The Lean formalisation of the proposition is
\texttt{continuous\_combiner\_bilinear\_classification} in
\texttt{Foundation.GeneralizedDAlembert}, conditional on the
\texttt{ContinuousCombinerAnalysisInputs} hypothesis package. The
strengthened Level-1 corollary,
\texttt{RCL\_form\_under\_continuous\_combiner}, is the form the
present paper revises into when the smoothness hypothesis is supplied;
under the assumption the polynomial caveat in the Level-1
\texttt{LogicAsFunctionalEquation} module drops automatically.

\begin{remark}[On the relation between the polynomial and continuous
versions]
The polynomial version of the main theorem (Theorem~\ref{thm:main}) and
the continuous-combiner version
(Proposition~\ref{prop:continuous-combiner-extension}) reach the same
combiner family on $\operatorname{Range}(F)^2$. The polynomial form has
the advantage of producing the bilinear family
$P(u, v) = 2u + 2v + c\,uv$ as an explicit polynomial on
$\mathbb{R}^2$; the continuous form has the advantage of dropping the
polynomial regularity. We use the polynomial version throughout the
body of the paper because it is the form proved in the recently
established and peer-reviewed Inevitability
Theorem~\cite{WashburnZlatanovicAllahyarov2026}.
\end{remark}

\subsection{The status of finite pairwise polynomial closure}
\label{subsec:conjecture}

\paragraph{The headline.} Finite pairwise polynomial closure is part
of the theorem's hypothesis. \emph{Without it, the theorem is false.}
The role of this subsection is to make that point sharply, exhibit the
witnesses, and explain why the resulting scope is the correct scope
for the identification.

\paragraph{Why it is essential.} The polynomial restriction (R2), now
named precisely as finite pairwise polynomial closure
(Definition~\ref{def:finite-pairwise-closure}), is not a redundant
regularity hypothesis. The counterexample of
Proposition~\ref{prop:finite-closure-necessary} shows that an operative
positive-ratio comparison can satisfy all of (L1)--(L4), (R1), (B1),
(B2), admit a continuous symmetric combining rule $\Phi$ that is closed
under iteration on $\operatorname{Range}(F)$, and yet $\Phi$ fail to be
polynomial of degree at most two: the example $C(x, y) = (\ln(x/y))^4$
produces $\Phi(a, b) = 2a + 2b + 12\sqrt{ab}$ on
$\operatorname{Range}(F) = [0, \infty)$. Closure under iteration alone
is therefore not sufficient to force polynomiality, and the
conclusion of Theorem~\ref{thm:rcl-direct} fails on this witness.

This refutes a stronger conjecture proposed in earlier drafts of this
paper, that closure under iteration of the combining rule on
$\operatorname{Range}(F)$ implies real-analyticity of $\Phi$ and hence
polynomial degree at most two. The quartic-log example is closed under
iteration on $[0, \infty)$, but its combining rule is not real-analytic
at the origin (the $\sqrt{ab}$ term has a branch-point character there).
The earlier conjecture is therefore false, and the paper records this
correction explicitly.

The right framing is the one given in
Remark~\ref{rem:r2-status-update}: finite pairwise polynomial closure
is the structural content that distinguishes the Recognition
Composition Law family from broader compositional algebras of
positive-ratio comparison. Without it, the conclusion of
Theorem~\ref{thm:rcl-direct} fails. The polynomial restriction is
therefore part of the structural meaning of \emph{finite pairwise
algebra} on this domain, not a regularity carried alongside the four
classical conditions for technical convenience.

\paragraph{Analyticity does not suffice.} A second possible repair is
to replace polynomiality by real-analyticity of the combiner at the
origin. That repair also fails. The failure is more instructive than a
technical caveat: the Recognition Composition Law family is not
invariant under analytic reparameterisation of the cost coordinate.

\begin{proposition}[Analytic reparameterisation counterexample]
\label{prop:analytic-reparam-counterexample}
There exists a real-analytic local change of cost coordinate for which
the induced combiner is real-analytic at the origin and satisfies the
d'Alembert form, but is not a polynomial of total degree at most two in
the reparameterised variables.
\end{proposition}

\begin{proof}
Let
\[
K(t)=\cosh t - 1.
\]
Then $K$ satisfies the standard RCL combiner:
\[
K(t+u)+K(t-u)=2K(t)+2K(u)+2K(t)K(u).
\]
Now change the cost coordinate analytically by
\[
f(s)=s+s^2,\qquad H(t)=f(K(t)).
\]
The function $f$ is analytic near $0$ and has analytic local inverse
\[
g(a)=\frac{-1+\sqrt{1+4a}}{2}.
\]
Writing $p=g(a)$ and $q=g(b)$, so that $a=p+p^2$ and $b=q+q^2$, the two
standard RCL branch values for $K(t+u)$ and $K(t-u)$ have the form
\[
A+R,\qquad A-R,
\]
where
\[
A=p+q+pq,\qquad R^2=p(2+p)q(2+q).
\]
The reparameterised combiner is therefore
\[
\Phi(a,b)=f(A+R)+f(A-R).
\]
For $f(s)=s+s^2$, the odd powers of $R$ cancel:
\[
\Phi(a,b)=2A+2A^2+2R^2,
\]
with $p=g(a)$ and $q=g(b)$. This is real-analytic near $(0,0)$.

Restrict to the diagonal $a=b$. Then $p=q=s$ and $a=s+s^2$. We compute
\[
A=2s+s^2,\qquad R^2=s^2(2+s)^2,
\]
and hence
\[
\Phi(a,a)=2A+2A^2+2R^2
         =4s+18s^2+16s^3+4s^4.
\]
If $\Phi$ were a degree-two RCL-family combiner with the boundary
condition $\Phi(a,0)=2a$, then on the diagonal it would have the form
\[
\Phi(a,a)=4a+c a^2.
\]
Substituting $a=s+s^2$ gives
\[
4a+c a^2
=4s+(4+c)s^2+2c s^3+c s^4.
\]
Matching the $s^2$ coefficient forces $c=14$, while the $s^3$ coefficient
would then be $28$, not the actual value $16$. Contradiction.
\end{proof}

Thus even real-analyticity at the origin does not force polynomial degree
at most two. Finite pairwise polynomial closure is not a proxy for
analytic regularity; it is the actual structural condition selecting the
RCL family in the chosen cost coordinate.

\paragraph{Closure under iteration as a structural condition.} The
closure-under-iteration condition is still a useful structural notion;
it merely does not by itself force analyticity. We retain its
formal definition for use in the Lean formalisation and in further
investigations.

\begin{definition}[Closure under iteration]
\label{def:closure-under-iteration}
A combining rule $\Phi : \mathbb{R}^2 \to \mathbb{R}$ is \emph{closed
under iteration on $\operatorname{Range}(F)$} if (a) $\Phi$ is jointly
continuous on $\operatorname{Range}(F) \times \operatorname{Range}(F)$,
and (b) for any $u, v \in \operatorname{Range}(F)$,
$\Phi(u, v) \in \operatorname{Range}(F)$. Equivalently, the dynamical
system $(u, v) \mapsto \Phi(u, v)$ on $\operatorname{Range}(F)$ is
well-defined and continuous.
\end{definition}

\begin{remark}[The Lean status, updated]
\label{rem:lean-status-update}
The Level-3 Lean module \texttt{Foundation.PolynomialityFromLogic} is
written conditionally on a class assumption
\texttt{IteratedAnalyticityHolds} that asserts closure under iteration
on $\operatorname{Range}(F)$ implies real-analyticity of $\Phi$ on
$\operatorname{Range}(F)^2$. The quartic-log counterexample shows that
the assumption is, in general, false: an instance of the assumption
fails for the explicit example $C(x, y) = (\ln(x/y))^4$. The Lean
formalisation therefore records the polynomiality conditional on a
class assumption that is not universally true. We retain the module as
a record of the structural consequences of closure under iteration
that hold without further regularity, but we mark
\texttt{IteratedAnalyticityHolds} as a hypothesis that is not
universally satisfied. The analytic reparameterisation counterexample
in Proposition~\ref{prop:analytic-reparam-counterexample}, formalised in
the module \texttt{Foundation.LogicAsFunctionalEquation.AnalyticCounterexample},
shows that even replacing \texttt{IteratedAnalyticityHolds} with
real-analyticity at the origin would not force polynomial degree at most
two. The correct Lean status is therefore: finite pairwise polynomial
closure is the theorem's hypothesis, not a removable regularity condition.
\end{remark}

\paragraph{Structural consequences of closure under iteration.} The
closure-under-iteration condition still has useful structural
consequences that hold without further regularity. The Lean module
\texttt{Foundation.PolynomialityFromLogic} establishes:
\begin{itemize}[leftmargin=*]
\item \emph{Diagonal continuity.} The diagonal map
$v \mapsto \Phi(v, v)$ is continuous on $\operatorname{Range}(F)$.
(Lean theorem \texttt{diagonal\_continuous\_on\_range}, fully proved.)
\item \emph{Iterated continuity.} For every $n \in \mathbb{N}$, there
exists a continuous map $\varphi_n :
\operatorname{Range}(F) \to \operatorname{Range}(F)$ representing the
$n$-fold iterate of $\Phi$ on the diagonal. (Lean theorem
\texttt{iterate\_continuous\_on\_range}, fully proved.)
\end{itemize}
These are direct consequences of joint continuity and closure on the
range, and they remain valid for the quartic-log counterexample as
well as for any RCL-family solution. They do not deliver analyticity,
and analyticity would not be enough anyway. Proposition~\ref{prop:analytic-reparam-counterexample}
shows that analytic changes of cost coordinate preserve analytic
compositionality while changing the polynomial degree of the combiner.

\paragraph{Final status of R2.} Finite pairwise polynomial closure
(R2) is final in this paper. It is the exact coordinate-level
condition under which the RCL family is forced. Dropping it gives
broader continuous-combiner algebras with admissible non-RCL
solutions; replacing it by real-analyticity gives analytically
reparameterised RCL algebras that are still not polynomial of degree
at most two. The theorem is sharp as stated, in both directions: it
is not the universal claim ``every continuous comparison operator
satisfying the four classical conditions yields the RCL family,''
which is false; it is the structurally narrower identification
``finite pairwise polynomial comparison semantics of classical
invariance imply the RCL family,'' which is true and is what
Theorem~\ref{thm:rcl-direct} proves. The two paragraphs of
\S\ref{subsec:scope-statement} state this scope in its sharpest
form; the rest of this paper is what is true inside that scope.

\subsection{The discrete and modal cases}

The present result is stated for continuous comparisons of positive ratios.
The discrete propositional case, where comparisons are between truth
values and the algebraic structure is Boolean rather than
multiplicative, is treated as a separate realization in the second
companion paper~\cite{Washburn2026UniversalForcing}. There the
universal-property characterisation of the iteration object handles
the parity collapse from the infinite iteration sequence onto the
two-element Boolean carrier, so the meta-theorem applies even though
the carrier image is finite. The classical Stone representation
theorem~\cite{Stone1936} is suggestive but not directly comparable to
the present rigidity, which lives at the cost-functional level rather
than at the algebra level. Modal, intuitionistic, and higher-order
extensions each have their own formal articulation of the laws of
logic, and the question of whether each admits a rigidity-style
identification with a unique mathematical object remains open and is
not addressed here.

\subsection{Beyond the bilinear branch}\label{subsec:additive-branch}

The case $c = 0$ in $(\mathrm{RCL})$ corresponds, after passage to
logarithmic coordinates, to a quadratic-logarithm cost
$F(x) = \tfrac{1}{2}(\ln x)^2$ under the unit log-curvature calibration.
This is also a canonical cost, distinct from $J$, parametrized by $c = 0$.
Whether the additive branch admits a separate logical interpretation, or
whether the Aristotelian conditions (perhaps together with an additional
structural condition) select the bilinear $c \neq 0$ branch over the
additive $c = 0$ branch, is a question we do not resolve here. See
\cite[Section~3]{WashburnRahnamai2026} for the optimization-theoretic
analysis of the relationship between the two branches.

\subsection{Numerical scope and dimension}

The present result is one-dimensional in the sense that the comparison
operator is defined on $\mathbb{R}_{>0} \times \mathbb{R}_{>0}$. The
multidimensional extension, comparisons of $n$-tuples of positive
quantities, is treated in
\cite[Section~4]{WashburnZlatanovicAllahyarov2026}, where it is shown that
for $c \neq 0$, every solution depends only on a single linear
combination of coordinate logarithms. Thus the present rigidity is
genuinely one-dimensional in its admissible solutions even when stated on
$\mathbb{R}_{>0}^n$. The dimensional rigidity of physical space, derived
in the parallel program~\cite{PardoGuerraThapaWashburnAllahyarov2026},
provides the operational interpretation of why three-dimensional
recognition geometry is the natural setting for physical applications;
the connection to the present result, on the level of laws of logic, is a
direction for future work.

\section{Lean Formalisation: Summary}\label{sec:lean-status}

The Translation Theorem (Theorem~\ref{thm:translation}) is formalised
in the canonical Lean~4 Recognition Science
library~\cite{LeanLibrary2026} as the module
\texttt{Foundation.LogicAsFunctionalEquation}, with no residual
\texttt{sorry}s, building against the
peer-reviewed~\cite{WashburnZlatanovicAllahyarov2026,
WashburnZlatanovic2026} modules. The continuous-combiner extension
(Proposition~\ref{prop:continuous-combiner-extension}) is formalised
in \texttt{Foundation.GeneralizedDAlembert} with no residual
\texttt{sorry}s and no top-level axioms. Two structural inputs that
were previously named axioms in this module have changed status. The
H-side classification (the case $\Phi(u, v) = 2 u v$ of the
generalised d'Alembert equation) is now a discharged theorem,
\texttt{aczel\_kannappan\_continuous\_dAlembert}, proved by reduction
to \texttt{Cost.FunctionalEquation.dAlembert\_classification}. The
combiner-side classification has been split: an explicit hypothesis
package, \texttt{ContinuousCombinerAnalysisInputs}, with three named
analytical components (mollifier finite-order smoothness, second
derivative identity, $\psi$-affine completion), now stands in for the
old monolithic \texttt{GeneralizedAczelSmoothnessPackage}. The
quartic-log counterexample is formalised in the same module and shows
that this hypothesis package is \emph{not} automatic from continuity:
the second-derivative identity fails for the operator
$C(x, y) = (\ln(x/y))^4$. The bilinear conclusion is therefore exposed
behind the explicit hypothesis package rather than as a continuity-only
consequence. The former \texttt{IteratedAnalyticityHolds} assumption has
been removed from \texttt{Foundation.PolynomialityFromLogic}, which
also has no residual \texttt{sorry}s. The Lean formalisation verifies
the mathematics of the translation given the encoded conditions; it
does not and cannot verify the faithfulness of the Aristotelian
labelling, which is a definitional choice defended in
Section~\ref{subsec:encoding-defenses}. Module-level details, the
discharge of \texttt{aczel\_kannappan\_continuous\_dAlembert}, the
contents of the \texttt{ContinuousCombinerAnalysisInputs} package, and
the build verification appear in Appendix~\ref{app:lean}.

\section{Conclusion}\label{sec:conclusion}

The paper proposes an operational encoding of the four
Aristotelian-style conditions on a continuous comparison operator on
positive ratios and carries two named regularity hypotheses (continuity
of the operator, polynomial degree-two closure of the combiner)
alongside them. The Translation Theorem (Theorem~\ref{thm:translation})
reduces the encoded conditions to the hypotheses of the d'Alembert
Inevitability Theorem~\cite{WashburnZlatanovicAllahyarov2026}, the
derived cost satisfies the Recognition Composition Law $F(xy) +
F(x/y) = 2F(x) + 2F(y) + c\,F(x)F(y)$, and under convexity,
non-negativity, and unit log-curvature calibration the family splits
into a bilinear branch with calibrated representative
$J(x) = \tfrac{1}{2}(x + x^{-1}) - 1$ and an additive branch with
calibrated representative $\tfrac{1}{2}(\ln x)^2$.

Three further results upgrade the rigidity into an identification.
The canonicality theorem (Theorem~\ref{thm:canonicality}) shows that
under the magnitude-of-mismatch interpretation, the encoding
(L1)--(L4) is the unique operator-level reading of the four classical
Aristotelian conditions. The reality theorem
(Theorem~\ref{thm:reality-satisfies-logic}) shows that any structure
with definite states, single-valued comparison, decidability, and
compositional inference automatically satisfies (L1)--(L4) on its
comparison operator. Combined with the rigidity machinery, these
two theorems yield Corollary~\ref{cor:rcl-logic-reality}: on the
operative domain, the recognition composition law, the law of logic,
and the structural form of any reality on this domain are the same
mathematical object, with calibrated representatives $J$ and
$\tfrac{1}{2}(\ln x)^2$. The operative domain is broad: every pairwise
physical measurement, every continuous likelihood ratio, every
continuous distance or divergence on a positive-quantity state space.

The most substantive boundary result is negative, and is the source
of the theorem's exact scope. Finite pairwise polynomial closure of
the combiner is not removable. The quartic-log counterexample of
Proposition~\ref{prop:finite-closure-necessary} exhibits an operative
comparison $C(x, y) = (\ln(x/y))^4$ whose continuous combiner
$\Phi(a, b) = 2a + 2b + 12\sqrt{ab}$ is closed under iteration on
$[0, \infty)$ but is not polynomial; the analytic-reparameterisation
counterexample of Proposition~\ref{prop:analytic-reparam-counterexample}
further shows that even real-analyticity at the origin is not enough.
Without finite pairwise polynomial closure the theorem is false. Read
correctly, this is not a weakness but a precise scope statement: the
identification proved here is not ``every comparison operator
satisfying the four classical conditions yields the Recognition
Composition Law,'' but the structurally narrower
\emph{finite pairwise polynomial comparison semantics of classical
invariance imply the Recognition Composition Law}. Polynomial
closure is what names the compositional class for which the
identification holds; outside that class, broader continuous-combiner
algebras are perfectly admissible alternatives, and they are exhibited
in the same paper. Whether other compositional classes admit their
own rigidity theorems is a separate open question.

Twenty-four centuries separate Aristotle's three laws of thought from
the present paper. In the interval, the laws have been formalised in
many vocabularies: Boole's algebra, Frege's calculus, Russell and
Whitehead's logicism, Hilbert's metalanguage, Tarski's semantics,
Gentzen's natural deduction, Curry-Howard, Lawvere's categorical
logic, Girard's linear logic, the univalent foundations program. Each
was a name for the laws in a chosen syntactic vocabulary. None
exhibited the laws as a single mathematical object whose form was
forced by their structural content. The present paper does so on the
operative domain.

\paragraph{Companion installments and remaining open questions.} The
arithmetic content of the encoding is recovered in the first
companion paper~\cite{Washburn2026Arithmetic} as the canonical orbit
of any non-trivial generator inside $(\mathbb{R}_{>0}, \cdot)$,
extended through the integers, rationals, reals, and complex carrier.
The realization-invariance of that arithmetic is the meta-theorem of
the second companion paper~\cite{Washburn2026UniversalForcing}. What
remains open are: the modal, intuitionistic, and higher-order
extensions of the encoding; whether a further structural condition
distinguishes $J$ from $\tfrac{1}{2}(\ln x)^2$ within the joint set
of (L1)--(L4) and (R1)--(R2); and the discharge of the explicit
hypothesis package \texttt{ContinuousCombinerAnalysisInputs} that
witnesses the obstruction recorded in
Section~\ref{sec:lean-status}.

Section~\ref{subsec:implications} discusses the specific bearings of
the result on the invented-versus-discovered debate, on the
mathematical-intuition / physical-reality question, on the
foundations of mathematics, on the physical-laws-as-conventions
question, on comparison operators across physical and cognitive
substrates, and on the classical identification of existence with
logical consistency.

\bigskip

\noindent\textbf{Funding.} This research received no external funding.

\noindent\textbf{Conflicts of interest.} The author declares no conflict of
interest.

\noindent\textbf{Data availability.} No new data were created or analysed
in this study.

\appendix

\section{Lean Formalisation: Detailed Status}\label{app:lean}

The contents of the present paper are formalised in the canonical
Lean~4 Recognition Science library~\cite{LeanLibrary2026}. This
appendix records the status of each new module, the precise residual
\texttt{sorry} count, the build verification, and the documented
proof strategies for the two named class assumptions on which the
conditional results depend.

\subsection{The three modules and their residuals}\label{app:modules}

\paragraph{Level~1: \texttt{Foundation.LogicAsFunctionalEquation}.}
The module formalises the encoding of the four Aristotelian conditions
(L1)--(L4), the bridge condition (B1), the non-triviality clause (B2),
the named regularity hypotheses (R1) and (R2), and the four short
translation lemmas connecting the encoded conditions to the
hypotheses of the d'Alembert Inevitability
Theorem~\cite{WashburnZlatanovicAllahyarov2026}. The main theorem
(\texttt{laws\_of\_logic\_imply\_dalembert\_hypotheses}) and the
Recognition Composition Law conclusion
(\texttt{RCL\_is\_unique\_functional\_form\_of\_logic}) are stated and
proved in full. The cost corollary
(\texttt{J\_is\_unique\_cost\_under\_logic}) invokes the
already-formalised \texttt{washburn\_uniqueness\_aczel} from
\texttt{Cost.FunctionalEquation} as a black box. The module contains
\textbf{0 \texttt{sorry}} and builds against the existing
peer-reviewed~\cite{WashburnZlatanovicAllahyarov2026,
WashburnZlatanovic2026} modules.

\paragraph{Level~2: \texttt{Foundation.GeneralizedDAlembert}.}
The module weakens the polynomial-degree-two regularity hypothesis
(R2) to continuity of the combiner alone. The structural results that
were already in place (boundary identities, symmetry transfer to the
closure of the range by density, the Acz\'el iteration substitution,
and the preconnected-image-implies-order-connected route to the
range-contains-an-interval property) are preserved. Two further
upgrades are now in place.

First, the H-side d'Alembert classification — the single deepest
classical input on which the original Level-2 conditional theorem was
keyed — is now a discharged theorem,
\texttt{aczel\_kannappan\_continuous\_dAlembert}. It states that any
continuous $H : \mathbb{R} \to \mathbb{R}$ with $H(0) = 1$ satisfying
$H(t + u) + H(t - u) = 2 H(t) H(u)$ is one of $1$, $\cosh(\alpha\,\cdot)$,
or $\cos(\alpha\,\cdot)$. The proof reduces to
\texttt{Cost.FunctionalEquation.dAlembert\_classification}, which is
itself built from the integration bootstrap, the universal
d'Alembert-to-ODE derivation, and the three ODE uniqueness lemmas
proved in \texttt{Cost/FunctionalEquation.lean}. No external classical
axiom enters the chain.

Second, the combiner-side classification has been refined. The old
monolithic class assumption \texttt{GeneralizedAczelSmoothnessPackage}
is replaced by an explicit three-part hypothesis package,
\texttt{ContinuousCombinerAnalysisInputs}, with named components
\texttt{finite\_smoothness} (mollifier-driven finite-order derivative
control on the log-coordinate cost),
\texttt{second\_derivative} (the
$2 G''(t) = \psi(G(t)) \cdot G''(0)$ identity), and
\texttt{psi\_affine} (Stetk{\ae}r-style affineness of $\psi$ on the
image of the log-cost). The bilinear classification
\texttt{continuous\_combiner\_bilinear\_classification} consumes this
package and assembles the bilinear conclusion through the discharged
H-side classification, the explicit log-bilinear lift
\texttt{log\_bilinear\_affine\_lift\_dAlembert}, and the
positive-ratio bilinear assembly
\texttt{log\_bilinear\_positive\_cost\_bilinear}. The same module now
formally records the obstruction: the quartic log-cost
$C(x, y) = (\ln(x/y))^4$ refutes any attempt to derive
\texttt{second\_derivative} from continuity alone. The package is
therefore exposed as a hypothesis surface, not as an automatic
consequence.

The module contains \textbf{0 \texttt{sorry}} and uses no top-level
\texttt{axiom}.

\paragraph{Level~3: \texttt{Foundation.PolynomialityFromLogic}.}
The module formalises the closure-under-iteration framing of the
strengthened (L4), proves that the algebraic and topological
consequences of closure-under-iteration follow without further
assumption (\texttt{diagonal\_continuous\_on\_range},
\texttt{iterate\_continuous\_on\_range}), and no longer states a
Level-3 main theorem. The former class assumption
\texttt{IteratedAnalyticityHolds} has been removed because it is false
in general, as shown by the quartic-log counterexample. The stronger
repair (assuming analyticity at the origin directly) is also false, as
shown by the analytic reparameterisation counterexample formalised in
\texttt{Foundation.LogicAsFunctionalEquation.AnalyticCounterexample}.
The module contains \textbf{0 \texttt{sorry}} after the cleanup.

\subsection{Build verification}\label{app:build}

The relevant modules build cleanly against the existing canonical
library. The build is verified by the standard \texttt{lake} commands
\texttt{lake build IndisputableMonolith.Foundation.LogicAsFunctionalEquation},
\texttt{lake build IndisputableMonolith.Foundation.GeneralizedDAlembert},
\texttt{lake build IndisputableMonolith.Foundation.PolynomialityFromLogic},
and the counterexample module
\texttt{lake build IndisputableMonolith.Foundation.LogicAsFunctionalEquation.AnalyticCounterexample}.
The whole-library build is reported in the current audit notes.

\subsection{Total residual sorry count}\label{app:sorry}

All four cited modules now contain \textbf{0 \texttt{sorry}}. The two
sorries previously present in \texttt{Foundation.GeneralizedDAlembert},
both keyed on the old \texttt{GeneralizedAczelSmoothnessPackage}, are
discharged: the H-side step is now the theorem
\texttt{aczel\_kannappan\_continuous\_dAlembert}, and the combiner-side
step is exposed as the explicit hypothesis package
\texttt{ContinuousCombinerAnalysisInputs} consumed by the bilinear
classification. The former \texttt{IteratedAnalyticityHolds} sorry was
removed along with the false class assumption. The main finite-pairwise
polynomial theorem and both counterexamples do not depend on the
hypothesis package.

\subsection{What the formalisation does and does not establish}\label{app:scope}

The formalisation establishes, to the standard of a Lean~4 proof
checked by the \texttt{lake} build, that the derivation chain from
the encoded Aristotelian-style conditions, through the d'Alembert
Inevitability Theorem, to the canonical reciprocal cost $J$ on the
bilinear branch is mathematically correct \emph{conditional on the
polynomial-degree-two regularity hypothesis}. The main translation
theorem of Level~1
(\texttt{laws\_of\_logic\_imply\_dalembert\_hypotheses}) and the
Recognition Composition Law conclusion
(\texttt{RCL\_is\_unique\_functional\_form\_of\_logic}) are
unconditional theorems in the Lean library, with zero \texttt{sorry}s
in the chain that produces them. Under the additional explicit
hypothesis package \texttt{ContinuousCombinerAnalysisInputs}, the
polynomial regularity is replaceable by continuity of the combiner
alone (Level~2); the same module records the quartic-log obstruction
that prevents the package from being derived from continuity alone.
There is no Level~3 replacement by closure under iteration: that route
is refuted by the counterexamples above.

A distinction worth flagging. The Lean formalisation verifies the
\emph{mathematics} of the translation: given the encoded conditions
as Lean definitions, the named conclusions follow by proofs the
machine has checked. The formalisation does \emph{not} verify the
\emph{faithfulness of the Aristotelian labelling} of those Lean
definitions, and cannot. The identification of $C(x, x) = 0$ with
identity, of $C(x, y) = C(y, x)$ with non-contradiction, of totality
with excluded middle, and of route-independence with composition
consistency is a definitional choice, defended in
Section~\ref{subsec:encoding-defenses}. A machine verifies theorems,
not definitions. A reader who rejects the encoding, or proposes an
alternative, derives different conclusions on the continuous-ratio
domain, which would have to be formalised separately. The Lean
library reports what the present encoding yields.

The formalisation does not establish that classical logic is
reducible to cost-physics or cost-physics to classical logic. It
establishes that on the continuous-ratio comparison domain, under
the encoding adopted here and finite pairwise polynomial closure, the
hypotheses required by both readings are the same hypotheses, with the
same family of canonical calibrated solutions (the $\alpha$-family on
the bilinear branch, with $J$ as its normalised representative; the
quadratic-logarithm form on the additive branch). The counterexamples
show that neither arbitrary continuous composition nor analytic
composition alone is enough to recover the polynomial closure.

\begin{thebibliography}{99}

\bibitem{Aristotle}
Aristotle.
\textit{Prior Analytics}.
Trans. R.~Smith, Hackett, Indianapolis, 1989. Original work c.~350~BCE.

\bibitem{LeanLibrary2026}
J.~Washburn et al.
\textit{The Recognition Science Lean~4 library.}
Canonical repository, 2026. Modules
\texttt{RecognitionScience.Foundation.LogicAsFunctionalEquation},
\texttt{RecognitionScience.Foundation.GeneralizedDAlembert},
\texttt{RecognitionScience.Foundation.PolynomialityFromLogic}, and
\texttt{RecognitionScience.Foundation.LogicAsFunctionalEquation.AnalyticCounterexample}
formalise the contents of the present paper.
The four modules build cleanly together
(\texttt{lake build IndisputableMonolith.Foundation.LogicAsFunctionalEquation
IndisputableMonolith.Foundation.GeneralizedDAlembert
IndisputableMonolith.Foundation.PolynomialityFromLogic
IndisputableMonolith.Foundation.LogicAsFunctionalEquation.AnalyticCounterexample},
7822 jobs, exit~0).

\bibitem{Aczel1966}
J.~Acz\'el.
\textit{Lectures on Functional Equations and Their Applications.}
Academic Press, New York, 1966.

\bibitem{Boole1854}
G.~Boole.
\textit{An Investigation of the Laws of Thought, on which are Founded the
Mathematical Theories of Logic and Probabilities.}
Walton and Maberly, London, 1854.

\bibitem{Curry1934}
H.~B.~Curry.
Functionality in combinatory logic.
\textit{Proc. Natl. Acad. Sci. USA} \textbf{20} (1934), 584--590.

\bibitem{Feynman1964}
R.~P.~Feynman.
\textit{The Character of Physical Law.}
Messenger Lectures, Cornell University, 1964; published MIT Press, 1965.

\bibitem{Frege1879}
G.~Frege.
\textit{Begriffsschrift, eine der arithmetischen nachgebildete Formelsprache
des reinen Denkens.}
Halle, 1879.

\bibitem{Gentzen1935}
G.~Gentzen.
Untersuchungen \"uber das logische Schliessen, I~\&~II.
\textit{Mathematische Zeitschrift} \textbf{39} (1935), 176--210, 405--431.

\bibitem{Girard1987}
J.-Y.~Girard.
Linear logic.
\textit{Theoretical Computer Science} \textbf{50} (1987), 1--101.

\bibitem{Goedel1931}
K.~G\"odel.
\"Uber formal unentscheidbare S\"atze der Principia Mathematica und
verwandter Systeme, I.
\textit{Monatshefte f\"ur Mathematik und Physik} \textbf{38} (1931), 173--198.

\bibitem{Hilbert1925}
D.~Hilbert.
\"Uber das Unendliche.
\textit{Mathematische Annalen} \textbf{95} (1926), 161--190.

\bibitem{HoTT2013}
The Univalent Foundations Program.
\textit{Homotopy Type Theory: Univalent Foundations of Mathematics.}
Institute for Advanced Study, Princeton, 2013.

\bibitem{Howard1969}
W.~A.~Howard.
The formulae-as-types notion of construction.
Manuscript 1969; published in
\textit{To H.~B.~Curry: Essays on Combinatory Logic, Lambda Calculus and
Formalism},
ed.~J.~P.~Seldin and J.~R.~Hindley, Academic Press, 1980, 479--490.

\bibitem{Kannappan2009}
P.~Kannappan.
\textit{Functional Equations and Inequalities with Applications.}
Springer, 2009.

\bibitem{Kuczma2009}
M.~Kuczma.
\textit{An Introduction to the Theory of Functional Equations and
Inequalities,} 2nd ed.
Birkh\"auser, 2009.

\bibitem{Lawvere1964}
F.~W.~Lawvere.
An elementary theory of the category of sets.
\textit{Proc. Natl. Acad. Sci. USA} \textbf{52} (1964), 1506--1511.

\bibitem{PardoGuerraThapaWashburnAllahyarov2026}
S.~Pardo-Guerra, A.~Thapa, J.~Washburn, and E.~Allahyarov.
On dimensional constraints in recognition geometry.
\textit{Mathematics} (MDPI), 2026, accepted.

\bibitem{RussellWhitehead1913}
B.~Russell and A.~N.~Whitehead.
\textit{Principia Mathematica,} vols.~I--III.
Cambridge University Press, 1910--1913.

\bibitem{Stetkaer2013}
H.~Stetk{\ae}r.
\textit{Functional Equations on Groups.}
World Scientific, 2013.

\bibitem{Stone1936}
M.~H.~Stone.
The theory of representations of Boolean algebras.
\textit{Trans. Amer. Math. Soc.} \textbf{40} (1936), 37--111.

\bibitem{Tarski1933}
A.~Tarski.
The concept of truth in formalized languages.
Original Polish 1933; English translation in
\textit{Logic, Semantics, Metamathematics,} Oxford, 1956.

\bibitem{WashburnRahnamai2026}
J.~Washburn and A.~Rahnamai Barghi.
Reciprocal convex costs for ratio matching: Axiomatic characterization.
\textit{Axioms} \textbf{15} (2026), 151.
\url{https://doi.org/10.3390/axioms15020151}.

\bibitem{WashburnZlatanovic2026}
J.~Washburn and M.~Zlatanovi\'c.
Uniqueness of the canonical reciprocal cost.
\textit{Mathematics} \textbf{14} (2026), 935.
\url{https://doi.org/10.3390/math14060935}.

\bibitem{WashburnZlatanovicAllahyarov2026}
J.~Washburn, M.~Zlatanovi\'c, and E.~Allahyarov.
The d'Alembert Inevitability Theorem.
\textit{Mathematics} \textbf{14}(8) (2026), 1386.
Submitted 17 March 2026; revised 13 April 2026; accepted 17 April 2026;
published 20 April 2026.
\url{https://doi.org/10.3390/math14081386}.

\bibitem{Washburn2026Arithmetic}
J.~Washburn.
\textit{Arithmetic from the Law of Logic, I:
The Continuous Positive-Ratio Realization,
with a Universal Forcing Program.}
Manuscript, April 2026.

\bibitem{Washburn2026UniversalForcing}
J.~Washburn.
\textit{The Universal Forcing Meta-Theorem:
One Law of Logic, One Arithmetic, Across All Realizations.}
Manuscript, April 2026.

\bibitem{Washburn2026FloorClosure}
J.~Washburn.
\textit{Absolute Floor Closure: Specifiability, Distinguishability, and
the Last Precondition of Recognition Science.}
Manuscript, April 2026.

\bibitem{Washburn2026Specifiability}
J.~Washburn.
\textit{Specifiability Forces Distinguishability.}
Manuscript, April 2026.

\bibitem{Washburn2026SelfBootstrap}
J.~Washburn.
\textit{The Self-Bootstrap of Distinguishability.}
Manuscript, April 2026.

\bibitem{Wigner1960}
E.~P.~Wigner.
The unreasonable effectiveness of mathematics in the natural sciences.
\textit{Comm. Pure Appl. Math.} \textbf{13} (1960), 1--14.

\end{thebibliography}

\end{document}
